\documentclass[11pt]{article}
\usepackage{amsmath,amssymb,amsthm}
\usepackage{hyperref}
\usepackage{microtype}
\usepackage{booktabs}
\emergencystretch=3em

% Remove end-of-proof square
\renewcommand{\qedsymbol}{}

% Notation
\newcommand{\sem}[1]{\left[\!\left[#1\right]\!\right]}
\newcommand{\Rel}{\mathcal{R}}
\newcommand{\Profiles}{\mathcal{Q}}
\newcommand{\Transitions}{\mathcal{T}}
\newcommand{\Lineages}{\mathcal{L}}
\newcommand{\Pred}{\mathsf{Pred}}
\newcommand{\exec}{\operatorname{exec}}
\newcommand{\ValidR}{\operatorname{Valid}_{R}}
\newcommand{\PoR}{\operatorname{PoR}}
\newcommand{\PoC}{\operatorname{PoC}}
\newcommand{\PoP}{\operatorname{PoP}}

\title{Beyond Proof of Possession:\\
A Relational Theory of Authority}
\author{Nicola Gallo}
\date{24 September 2026}

\newtheorem{definition}{Definition}
\newtheorem{lemma}{Lemma}
\newtheorem{theorem}{Theorem}
\newtheorem{proposition}{Proposition}
\newtheorem{corollary}{Corollary}
\newtheorem{remark}{Remark}
\newtheorem{assumption}{Assumption}

\begin{document}
\maketitle

\begin{abstract}
Proof-of-Possession (PoP) establishes control of an artifact such as a key,
token, credential, or capability. Previous work on Provenance Identity
Continuity (PIC) and Proof-of-Continuity (PoC) showed that authority can instead
be propagated through a discrete causal execution by requiring, at every hop,
a verifiable Proof of Relationship (PoR) to the predecessor together with
non-expansive authority propagation. In that model, authority decreases
monotonically along a causal chain, so an in-lineage exercise beyond the origin
context cannot be represented as valid behavior. The effect-level
confused-deputy result developed here additionally assumes complete mediation
and causal attribution.

This paper develops the relational consequence of that result. We distinguish
execution steps from the executors that realize them and classify executors by
authorization-relevant equivalence relations. Each equivalence relation induces
a quotient space in which physically different executors are identified when
they are indistinguishable with respect to the relevant relational
characteristic. Multiple quotient views may coexist for the same executor.
Predicates over these relational views may be composed with conjunction and
disjunction, producing a distributive lattice ordered by logical implication.

This exposes a second monotonicity in continuous authorization. Authority
contexts attenuate by set inclusion, while relational admissibility attenuates
by predicate refinement. Both are indexed by the same causal chain established
by PoR. Causal history grows by prefix extension, while authority and relational
freedom can only decrease. We define a product authorization state containing
both dimensions and prove dual monotonicity and relational closure: no valid
descendant can introduce either a privilege or a relationally admissible
transition excluded by the origin context.

Finally, we separate possession from authorization. Proof-of-Possession can
contribute holder-binding or authentication evidence. Control of an artifact
alone does not establish an executor's profile attributes: that step requires a
trusted binding from the artifact or credential to the executor and its
classes. Authorization consumes authenticated profile facts together with PoR,
the authority context \(C_i\), and the relational envelope \(\rho_i\). We prove
that the executor-dependent decision factors through the relational profile, so
a sound mechanism that establishes the same profile does not change that
decision. In the model developed here, successor-side PoP is therefore one
possible evidence mechanism, not a necessary premise of the authorization
semantics. We further exhibit a boundary-witnessed construction in which the
successor proves possession of nothing, conditional on explicit Witness
Assumptions; in that checkpoint mechanism, the cryptographic possession resides
in the signing keys of the boundary witnesses that attest the relationship.
The construction does not independently prove those environmental assumptions.
\end{abstract}

\section{Introduction}

Most deployed authorization mechanisms attach authority to something that an
actor possesses or controls: a bearer token, a sender-constrained token, a key,
a credential, a capability, or an ambient privilege. Such mechanisms are often
effective local authorization tools. Their common abstraction, however, is
spatial: the decision is made from authority that is available to the actor at
the point of use.

Proof-of-Continuity~\cite{gallo2026poc} introduced a different axis. A later execution step is valid
only if it is causally related to its predecessor and carries no more authority
than the predecessor received. Formally, for a causal chain
\[
s_0 \rightarrow s_1 \rightarrow \cdots \rightarrow s_n,
\]
each valid hop satisfies
\[
\PoR(s_i,s_{i+1})
\qquad\text{and}\qquad
C_{i+1}\subseteq C_i,
\]
where \(C_i\subseteq O\times R\) is the authority context carried at step
\(s_i\)~\cite{gallo2026poc}. This gives
\[
C_n\subseteq C_{n-1}\subseteq\cdots\subseteq C_0
\]
and therefore origin-bounded authority: a privilege absent from \(C_0\) cannot
appear later in the same valid lineage.

That result leaves a further question open. If PoR is explicitly a
\emph{relationship} between execution occurrences, what mathematical structure
does that relationship carry? More specifically, can authorization be expressed
over relational characteristics of the execution rather than over possession by
a particular holder?

This paper answers that question by introducing a relational space for
executors. Let \(X\) be the set of executor instances. An
authorization-relevant equivalence relation \(\sim_j\) partitions \(X\) into
classes of executors that are indistinguishable with respect to one relational
characteristic. The quotient
\[
\Rel_j := X/{\sim_j}
\]
therefore forgets all executor differences irrelevant to that characteristic.
Different relations induce different quotient views, so the same executor may
belong to different equivalence classes under different relations.

We then define relational predicates over combinations of these quotient views.
Predicates compose with logical operators such as \(\land\) and \(\lor\), and
logical implication induces an attenuation order:
\[
\rho' \preceq_R \rho
\quad\Longleftrightarrow\quad
\rho' \Rightarrow \rho.
\]
Thus a successor may preserve or restrict a relational authorization envelope,
but may not expand it.

The result is a dual monotonicity. Along one causal chain,
\[
C_{i+1}\subseteq C_i
\]
restricts authority, while
\[
\rho_{i+1}\Rightarrow\rho_i
\]
restricts relational admissibility. The lineage prefix itself moves in the
opposite direction: causal history accumulates. Hence continuity has a simple
structural form:
\[
\text{history grows, while authority and relational freedom shrink.}
\]

The paper makes six contributions:
\begin{enumerate}
    \item it separates execution steps from executors and gives executor
    relations a quotient-space semantics;
    \item it defines multiple relational views of the same executor and a
    combined relational profile space;
    \item it defines a positive relational predicate lattice with refinement
    ordered by implication;
    \item it proves dual monotonicity and relational closure along a PoR-linked
    causal chain;
    \item it shows that the lineage-forgetting projection of the previous
    paper is itself a relational quotient, so that its impossibility result for
    possession-based policies is an instance of the factorization property
    developed here; and
    \item it proves that the executor-dependent component of authorization
    factors through the relational profile, formalizes invariance under sound
    profile-evidence mechanisms that produce the same profile, and gives,
    relative to explicit Witness Assumptions, a construction in which neither
    profile nor PoR evidence requires successor-side possession.
\end{enumerate}

Monotone attenuation along a chain is not itself new; capability tokens and
recent delegation-chain proposals for AI agents already provide it
(Section~\ref{sec:related}). The novelty claimed here is monotone attenuation of
an \emph{executor-relational admissibility space}, composed with authority
attenuation and causal continuity. Section~\ref{sec:picx} maps a restricted
subset of the premises onto an open-source implementation and states the gaps
explicitly.

The result does not claim that PoP is useless, nor that deployed
possession-based authorization systems do not authorize. Rather, it separates
two roles that are often collapsed: \emph{evidence that an actor controls
something} and \emph{the semantic condition that makes a continuation
authorized}. In the model developed here, possession may help establish the
former. Authentication fixes executor/profile facts; PoR establishes causal
continuation; and \(C_i\) together with \(\rho_i\) bounds the authorization
envelope.

The picture is one of two layers. At admission, a sound evidence mechanism
establishes facts about the executor and fixes its profile \(q_J(x)\). PoP may
provide holder-binding evidence only when combined with a trusted binding from
the possessed artifact to the executor and the asserted classes; the evidence
mechanism need not use successor-side PoP. Once the profile is established,
PoR determines whether the occurrence is a causal continuation, \(C_i\) bounds
which authority may continue, and \(\rho_i\) both tests and bounds the executor
relationships through which it may continue.
Treating possession as the ground of authority collapses the two layers. The
collapse is harmless at a single hop, where capabilities work precisely because
designation and authority coincide. It becomes visible across hops, because a
possession-based decision is invariant under the lineage
(Remark~\ref{rem:instance}).

\section{Background and Previous Result}

\subsection{Authority as a causal continuation}

Let \(O\) be a set of operations and \(R\) a set of resources. A privilege is
\[
(o,r)\in O\times R.
\]
A discrete execution is a causal chain
\[
\pi=\langle (s_0,C_0),(s_1,C_1),\ldots,(s_n,C_n)\rangle,
\]
where \(s_i\) is an execution step and
\[
C_i\subseteq O\times R
\]
is its authority context.

The previous PoC model defines a single-hop Proof of Relationship by
\[
\PoR(s_i,s_{i+1}),
\]
which holds when \(s_{i+1}\) is a valid causal continuation of \(s_i\) within
the same execution lineage~\cite{gallo2026poc}. A transition is valid when it
is both causally related and non-expansive:
\[
\operatorname{Valid}(s_i,C_i,s_{i+1},C_{i+1})
\Longleftrightarrow
\PoR(s_i,s_{i+1})
\land
C_{i+1}\subseteq C_i.
\]

The whole chain satisfies Proof-of-Continuity when every adjacent transition is
valid:
\[
\PoC(\pi)
\Longleftrightarrow
\bigwedge_{i=0}^{n-1}
\left(
\PoR(s_i,s_{i+1})
\land
C_{i+1}\subseteq C_i
\right).
\]

\begin{proposition}[Origin-bounded authority, previous result]
If \(\PoC(\pi)\) holds, then
\[
C_n\subseteq C_{n-1}\subseteq\cdots\subseteq C_0.
\]
Consequently, if \((o,r)\notin C_0\), then
\[
(o,r)\notin C_i
\]
for every \(i\).
\end{proposition}

\begin{proof}
Immediate by transitivity of set inclusion.
\end{proof}

\subsection{Confused deputy as authority/causality mismatch}

The classical confused deputy occurs when a request causes a deputy to exercise
authority that does not belong to the request's authority context. Capabilities
address the classical single-invocation form by fusing designation and
authority~\cite{hardy1988confused,miller2003capability}. In a distributed or
multi-step execution, however, an executor may hold several unrelated authority
sources simultaneously.

Let a request originate with authority context \(C_0\). If an executor later
uses a privilege
\[
(o,r)\notin C_0
\]
because it independently possesses that privilege, then authority and causality
have diverged. The previous PoC result rules this out as valid behavior because
every valid descendant context is a subset of \(C_0\).

The present work preserves this result and adds a second closure dimension: a
valid descendant may not expand the relational class of executions through
which origin authority is permitted to continue.

\section{Execution Steps and Executors}

The previous model uses \(s_i\) for a formal execution step. Operationally,
that step is realized by an executor such as a service, workload, process,
function, tool, agent, or hardware component.

\begin{definition}[Executor space]
Let
\[
S
\]
be the set of execution steps and let
\[
X
\]
be the set of executor instances. Define
\[
\exec:S\rightarrow X.
\]
For each \(s_i\in S\), write
\[
x_i:=\exec(s_i).
\]
\end{definition}

A causal chain
\[
s_0\rightarrow s_1\rightarrow\cdots\rightarrow s_n
\]
therefore induces an executor sequence
\[
x_0,x_1,\ldots,x_n.
\]

The distinction between step and executor is essential. The same executor may
realize several distinct occurrences:
\[
s_i\neq s_j
\qquad\text{while}\qquad
x_i=x_j.
\]
Conversely, a lineage may cross many different executors.

Authorization therefore cannot be reduced to executor identity alone. The
relevant object is an \emph{occurrence} of an executor in a causal execution and
the relationships under which that occurrence is admitted.

\section{Relational Equivalence and Executor Quotients}

\subsection{Equivalence relative to an authorization characteristic}

We now formalize the idea that two physically different executors may be
equivalent for one authorization purpose.

\begin{definition}[Authorization-relevant equivalence]
Let \(\sim_j\) be an equivalence relation on \(X\):
\[
\sim_j\;\subseteq X\times X.
\]
For \(x,y\in X\),
\[
x\sim_j y
\]
means that \(x\) and \(y\) are indistinguishable with respect to relational
characteristic \(j\).
\end{definition}

Since \(\sim_j\) is an equivalence relation, it is reflexive, symmetric, and
transitive. It partitions \(X\) into disjoint equivalence classes.

\begin{definition}[Equivalence class]
For \(x\in X\),
\[
[x]_{\sim_j}
:=
\{y\in X:y\sim_j x\}.
\]
\end{definition}

The class contains every executor that is indistinguishable from \(x\) under
the selected relation.

\begin{definition}[Relational quotient]
The quotient induced by \(\sim_j\) is
\[
\boxed{
\Rel_j := X/{\sim_j}
}
\]
with
\[
\Rel_j
=
\{
[x]_{\sim_j}:x\in X
\}.
\]
We call \(\Rel_j\) a \emph{relational quotient}.
\end{definition}

The quotient operation intentionally discards distinctions that the relation
declares irrelevant. It does \emph{not} assert that two concrete executors are
identical.

\subsection{A concrete quotient example}

Let
\[
X=\{a,b,c,d,e,f\}.
\]
Suppose \(\sim_1\) induces the partition
\[
\{a,b\},
\qquad
\{c,d,e\},
\qquad
\{f\}.
\]
Then
\[
\Rel_1
=
X/{\sim_1}
=
\big\{
\{a,b\},
\{c,d,e\},
\{f\}
\big\}.
\]

Since
\[
a\sim_1 b,
\]
we have
\[
[a]_{\sim_1}
=
[b]_{\sim_1}
=
\{a,b\}.
\]

Thus the quotient identifies \(a\) and \(b\) only with respect to the
characteristic represented by \(\sim_1\).

Now let a second relation \(\sim_2\) induce
\[
\{a,c,f\},
\qquad
\{b,d\},
\qquad
\{e\}.
\]
Then
\[
\Rel_2
=
X/{\sim_2}
=
\big\{
\{a,c,f\},
\{b,d\},
\{e\}
\big\}.
\]

The same executor \(a\) therefore has
\[
[a]_{\sim_1}
=
\{a,b\}
\]
but
\[
[a]_{\sim_2}
=
\{a,c,f\}.
\]

There is no contradiction. The same executor is being observed through two
different authorization-relevant relations.

\subsection{Quotient maps}

Each equivalence relation has a canonical quotient map
\[
q_j:X\rightarrow \Rel_j
\]
defined by
\[
q_j(x)=[x]_{\sim_j}.
\]

The map says: retain only the relational class of \(x\) and forget its
irrelevant concrete differences.

This gives the intended reading:
\[
X/{\sim_j}
=
\text{``\(X\) viewed at the resolution determined by \(\sim_j\)''.}
\]

\section{Multiple Relational Views}

A single executor may participate in several authorization-relevant
relationships. Let \(J\) index a family
\[
\{\sim_j\}_{j\in J}.
\]

Each relation produces its own quotient:
\[
\Rel_j=X/{\sim_j}.
\]

\begin{definition}[Relational profile map]
Define
\[
q_J:X\rightarrow \prod_{j\in J}\Rel_j
\]
by
\[
q_J(x)
=
\big(
[x]_{\sim_j}
\big)_{j\in J}.
\]
The realizable relational profile space is
\[
\Profiles_J
:=
q_J(X)
\subseteq
\prod_{j\in J}\Rel_j.
\]
\end{definition}

The use of the image \(q_J(X)\), rather than the full product, records only
combinations of quotient classes realized by some executor.

\subsection{Combined equivalence}

The family of relations also induces a combined equivalence:
\[
x\equiv_J y
\quad\Longleftrightarrow\quad
\forall j\in J,\;x\sim_j y.
\]
Equivalently,
\[
\equiv_J
=
\mathop{\cap}\limits_{j\in J}\sim_j.
\]

Two executors are therefore equivalent under \(J\) precisely when their entire
relational profiles agree:
\[
x\equiv_J y
\quad\Longleftrightarrow\quad
q_J(x)=q_J(y).
\]

\begin{proposition}[Factorization through a quotient]\label{prop:factorization}
Let \(Y\) be a set, \(\equiv\) an equivalence relation on \(Y\), and
\(q:Y\rightarrow Y/{\equiv}\) the canonical quotient map. A policy
\[
A:Y\rightarrow\{0,1\}
\]
is invariant under \(\equiv\), i.e.
\[
y\equiv y'
\Longrightarrow
A(y)=A(y'),
\]
if and only if there exists
\[
\bar A:Y/{\equiv}\rightarrow\{0,1\}
\]
such that
\[
A=\bar A\circ q.
\]
When it exists, \(\bar A\) is unique.
\end{proposition}

\begin{proof}
If \(A=\bar A\circ q\) and \(y\equiv y'\), then \(q(y)=q(y')\), hence
\(A(y)=A(y')\). Conversely, for a class \(c\in Y/{\equiv}\) choose any \(y\)
with \(q(y)=c\) and define \(\bar A(c):=A(y)\). Invariance makes this value
independent of the chosen representative, and \(A=\bar A\circ q\) by
construction. Uniqueness follows because \(q\) is surjective.
\end{proof}

For executors, take \(Y=X\) and \(\equiv\;=\;\equiv_J\). Since
\(x\equiv_J y\) iff \(q_J(x)=q_J(y)\), the profile map induces a bijection
\(X/{\equiv_J}\cong\Profiles_J\). A policy on executors is therefore invariant
under \(\equiv_J\) iff it factors as
\[
A=\bar A\circ q_J,
\qquad
\bar A:\Profiles_J\rightarrow\{0,1\}.
\]

This is the formal statement that an authorization rule may operate on
relational equivalence classes rather than on concrete executor identity. The
proposition is the standard universal property of a quotient; its interest here
lies in what it unifies.

\begin{remark}[The previous result as a quotient instance]\label{rem:instance}
In~\cite{gallo2026poc}, an event is an occurrence
\(e=(o,r,\ell)\in E:=O\times R\times\Lineages\), and the projection
\(p:E\rightarrow O\times R\) forgets the lineage. Let \(\sim_{(o,r)}\) be the
kernel relation of \(p\),
\[
e\sim_{(o,r)}e'
\quad\Longleftrightarrow\quad
p(e)=p(e'),
\]
and \(q:E\rightarrow E/{\sim_{(o,r)}}\) its canonical quotient map. Since \(p\)
is surjective, there is a canonical bijection
\(\phi:E/{\sim_{(o,r)}}\rightarrow O\times R\) with \(p=\phi\circ q\). Up to
this bijection, \(p\) is the quotient map, and a policy factors through \(p\)
iff it factors through \(q\).
The possession-based policies defined there, which factor through \(p\) or
equivalently are lineage-invariant (Definition~7 of~\cite{gallo2026poc}), are
exactly the policies invariant under \(\sim_{(o,r)}\). The result that
lineage-invariant policies cannot individuate occurrences (Theorem~3
of~\cite{gallo2026poc}) is Proposition~\ref{prop:factorization} applied to
\(Y=E\) and \(\equiv\;=\;\sim_{(o,r)}\): two events in the same class receive the
same verdict. The companion result that any resolution reintroduces continuity
(Theorem~4 of~\cite{gallo2026poc}) is its contrapositive: a policy that
separates two events of one class is not invariant under \(\sim_{(o,r)}\), and
must read a finer relation, namely one that consults the lineage.

The present construction therefore does not add a structure beside that of the
previous paper; it contains it. The lineage-forgetting projection is one
relational quotient, and the executor relations \(\sim_j\) are further ones,
all governed by the same factorization property.
\end{remark}

\section{Relational Predicates}

The quotient spaces classify executors. Authorization additionally needs a
language for stating which relational profiles, and which transitions between
profiles, are admissible.

The two layers have different jobs. An equivalence relation \(\sim_j\) must be
reflexive, symmetric and transitive, and many security relations are none of
these: delegation, acts-for, trust, parent and child, supervision, causal
precedence. The quotients are therefore \emph{classification spaces}, not the
authorization relation itself. The authorization-relevant relation between a
predecessor and its successor is a predicate on
\(\Profiles_J\times\Profiles_J\), which may be directional and need not be an
equivalence. A supervisory relation, for instance, is an atom
\(\mathrm{Sup}(u,v)\) that holds when the class of \(u\) supervises the class
of \(v\).

\subsection{Transition profile space}

A causal hop involves two executors:
\[
x_i=\exec(s_i),
\qquad
x_{i+1}=\exec(s_{i+1}).
\]

Define the transition profile space
\[
\Transitions_J
:=
\Profiles_J\times\Profiles_J.
\]

The relational observation of a concrete hop is
\[
\tau_i
:=
\big(
q_J(x_i),
q_J(x_{i+1})
\big)
\in\Transitions_J.
\]

\subsection{Atomic and composite predicates}

\begin{definition}[Relational predicate]
A relational predicate is a function
\[
\rho:\Transitions_J\rightarrow\{0,1\}.
\]
Its semantic extension is
\[
\sem{\rho}
:=
\{\tau\in\Transitions_J:\rho(\tau)=1\}.
\]
\end{definition}

Let
\[
R_1,R_2,\ldots
\]
be atomic predicates over \(\Transitions_J\). We use the positive relational
language generated by
\[
\rho ::= R_k \mid (\rho\land\rho)\mid(\rho\lor\rho).
\]

We write \(\Pred(\Transitions_J)\) for the set of positive relational
predicates generated by this grammar from the atoms. Relational envelopes range
over \(\Pred(\Transitions_J)\), not over arbitrary functions
\(\Transitions_J\rightarrow\{0,1\}\).

\paragraph{Why no negation.}\label{sec:relpred}
The answer depends on what a boundary knows about an atom. Each atom is a fixed
subset of \(\Transitions_J\), and the profiles in \(\tau\) are verified. Some
atoms are decidable from \(\tau\) itself, such as equality or inequality of two
profile components. For those, the atom and its negation are equally
decidable, and nothing is lost by admitting \(\lnot\). The separation-of-duty
atom of Section~\ref{sec:sod} is of this kind, which is why it may contain
\(\notin\) and \(\neq\). For other atoms the extension is not known to the
boundary: a supervision relation that depends on an external organizational
chart, or a relation recorded in a delegation registry. The boundary then
learns membership only through evidence, and the truth of the atom must be kept
apart from the availability of evidence for it.

Let \(E\) be the evidence available to a boundary, and write
\(E\vdash R_k(\tau)\) when \(E\) supports \(R_k(\tau)\). The relation lifts to
positive predicates:
\[
E\vdash(\rho\land\sigma)(\tau)
\;\Longleftrightarrow\;
E\vdash\rho(\tau)\text{ and }E\vdash\sigma(\tau),
\]
\[
E\vdash(\rho\lor\sigma)(\tau)
\;\Longleftrightarrow\;
E\vdash\rho(\tau)\text{ or }E\vdash\sigma(\tau).
\]
We assume two properties of atoms. Evidence is \emph{sound} when
\(E\vdash R_k(\tau)\) implies \(R_k(\tau)\). It is \emph{monotone} when
\(E\subseteq E'\) and \(E\vdash R_k(\tau)\) imply \(E'\vdash R_k(\tau)\). Both
properties lift by induction to every positive predicate. By soundness, a
boundary that admits on \(E\vdash\rho_i(\tau_i)\) admits only transitions that
satisfy \(\rho_i\). By monotonicity, withholding evidence cannot produce an
admission that full evidence would deny. Negation breaks this for atoms of the
second kind: supporting \(\lnot R_k(\tau)\) requires evidence of absence, which
an open world does not provide, and reading ``not supported'' as ``false''
would let missing evidence cause admission.

The positive fragment is therefore the conservative choice, safe for both kinds
of atom. Negation is not excluded on mathematical grounds. With \(\lnot\)
added, the implication order and every attenuation argument below are
unchanged, and the lattice of Proposition~\ref{prop:lattice} becomes a Boolean
algebra. A deployment may admit \(\lnot\) on atoms decidable at the boundary, or
under a closed-world assumption. The validity conditions below are stated at the
level of truth; a boundary discharges them through sound evidence.

Conjunction and disjunction have their ordinary set semantics:
\[
\sem{\rho\land\sigma}
=
\sem{\rho}\cap\sem{\sigma},
\]
and
\[
\sem{\rho\lor\sigma}
=
\sem{\rho}\cup\sem{\sigma}.
\]

Thus one causal execution may require, for example,
\[
R_1\land R_2,
\]
another may allow
\[
R_1\lor R_2,
\]
and a later continuation may refine the condition to
\[
(R_1\lor R_2)\land R_3.
\]

\section{The Relational Predicate Order}

We now define the order that will carry relational continuity.

\begin{definition}[Relational refinement]
For relational predicates \(\rho'\) and \(\rho\), define
\[
\rho'\preceq_R\rho
\]
iff
\[
\rho'\Rightarrow\rho,
\]
that is, \(\rho'(\tau)=1\) implies \(\rho(\tau)=1\) for every
\(\tau\in\Transitions_J\). Equivalently,
\[
\sem{\rho'}
\subseteq
\sem{\rho}.
\]
\end{definition}

Throughout, \(\Rightarrow\) between predicates denotes this extensional
implication, under the fixed interpretation of the atoms.

A stronger predicate therefore denotes a smaller set of admissible transitions.

For example,
\[
R_1
\Rightarrow
R_1\lor R_2,
\]
so
\[
R_1
\preceq_R
R_1\lor R_2.
\]

Likewise,
\[
R_1\land R_3
\Rightarrow
R_1,
\]
so
\[
R_1\land R_3
\preceq_R
R_1.
\]

Hence the sequence
\[
R_1\lor R_2
\;\succeq_R\;
R_1
\;\succeq_R\;
R_1\land R_3
\]
is monotonically restrictive.

\begin{proposition}[Relational predicate lattice]\label{prop:lattice}
Identify positive relational predicates over \(\Transitions_J\) that have the
same semantic extension. Ordered by \(\preceq_R\), the resulting classes form a
distributive lattice with meet \(\land\) and join \(\lor\).
\end{proposition}

\begin{proof}
Under the semantic map \(\rho\mapsto\sem{\rho}\), conjunction maps to set
intersection and disjunction maps to set union. The extensions of positive
predicates are therefore closed under \(\cap\) and \(\cup\), and form a
sublattice of the powerset of \(\Transitions_J\) ordered by inclusion. Every
sublattice of a distributive lattice is distributive.
\end{proof}

A second identification is possible: the free positive algebra, in which
formulas are identified when they are equivalent under \emph{every}
interpretation of the atoms. It is finer than the extensional one, because two
atoms may have the same extension without being propositionally equivalent. The
closure results below use only the extensional order. The syntactic check that
follows is the interpretation-independent one.

This lattice is the formal version of a relational authorization space. A
successor may move downward by refinement, thereby accepting fewer relational
transitions, but may not move upward without an explicit re-authorization or a
new origin.

\paragraph{Checking refinement.}\label{sec:refinement}
As a semantic condition, \(\rho'\Rightarrow\rho\) depends on the extensions of
the atoms. A verifier need not compute them, because an
interpretation-independent refinement can be checked syntactically. Every
formula in the positive grammar is finite and can be written in disjunctive
normal form. Regard each conjunction as the finite set of atoms occurring in
it.

\begin{lemma}[Positive-DNF implication criterion]\label{lem:dnf}
Let
\[
\rho'=\bigvee_{a\in A}\bigwedge_{p\in P_a}p,
\qquad
\rho=\bigvee_{b\in B}\bigwedge_{p\in Q_b}p
\]
be positive formulas in DNF. Then \(\rho'\Rightarrow\rho\) holds under every
Boolean interpretation of the atoms iff
\[
\forall a\in A\;\exists b\in B:\quad Q_b\subseteq P_a.
\]
Equivalently, every disjunct of \(\rho'\) contains all atoms of some disjunct
of \(\rho\).
\end{lemma}

\begin{proof}
If \(Q_b\subseteq P_a\), every valuation satisfying the conjunction indexed by
\(P_a\) also satisfies the one indexed by \(Q_b\). The displayed condition
therefore makes every disjunct of \(\rho'\) imply \(\rho\), and hence makes
\(\rho'\Rightarrow\rho\) valid under every interpretation.

Conversely, suppose that for some \(a\in A\), no \(Q_b\) is contained in
\(P_a\). Assign true exactly to the atoms in \(P_a\) and false to every other
atom occurring in either formula. This valuation satisfies the \(a\)-th
disjunct of \(\rho'\). For every \(b\in B\), choose
\(p_b\in Q_b\setminus P_a\); that atom is false, so no disjunct of \(\rho\) is
satisfied. Thus \(\rho'\not\Rightarrow\rho\), proving the contrapositive.
\end{proof}

In particular, conjoining a further predicate
(\(\rho'=\rho\land\sigma\)) and deleting disjuncts are always refinements. The
criterion of Lemma~\ref{lem:dnf} is sound and decidable from the formulas alone.
It is conservative for the fixed semantic interpretation used by the model:
particular atom extensions may make \(\rho'\Rightarrow\rho\) hold even when the
interpretation-independent criterion fails.

\section{Relational Continuity Along a Causal Chain}

\subsection{Relational envelopes}

Associate with every execution step \(s_i\) both:

\begin{enumerate}
    \item an authority context
    \[
    C_i\subseteq O\times R,
    \]
    and
    \item a relational envelope
    \[
    \rho_i\in\Pred(\Transitions_J).
    \]
\end{enumerate}

The relational envelope specifies the class of executor-profile transitions
that remain admissible at that point in the lineage.

\subsection{Extended valid transition}

\begin{definition}[Relationally valid transition]
Let
\[
x_i=\exec(s_i),
\qquad
x_{i+1}=\exec(s_{i+1}),
\]
and
\[
\tau_i=
\big(
q_J(x_i),
q_J(x_{i+1})
\big).
\]
The transition from
\[
(s_i,C_i,\rho_i)
\]
to
\[
(s_{i+1},C_{i+1},\rho_{i+1})
\]
is relationally valid iff
\[
\begin{aligned}
\ValidR
(s_i,C_i,\rho_i,s_{i+1},C_{i+1},\rho_{i+1})
\Longleftrightarrow\;&
\PoR(s_i,s_{i+1})
\\
&\land\; C_{i+1}\subseteq C_i
\\
&\land\; \rho_i(\tau_i)
\\
&\land\;(\rho_{i+1}\Rightarrow\rho_i).
\end{aligned}
\]
\end{definition}

The four conjuncts have different roles:

\begin{enumerate}
    \item \(\PoR(s_i,s_{i+1})\) establishes causal continuation;
    \item \(C_{i+1}\subseteq C_i\) prevents authority expansion;
    \item \(\rho_i(\tau_i)\) requires the concrete executor transition to
    satisfy the currently admitted relation; and
    \item \(\rho_{i+1}\Rightarrow\rho_i\) prevents relational expansion.
\end{enumerate}

Two design points concern the third conjunct. First, hop \(i\) is governed by
the predecessor's envelope \(\rho_i\), not by the successor's \(\rho_{i+1}\).
The envelope carried at \(s_i\) constrains the transition that continues
\(s_i\); the envelope \(\rho_{i+1}\) constrains the transitions that will
continue \(s_{i+1}\). Requiring \(\rho_{i+1}(\tau_i)\) as well would impose a
stronger, retroactive policy that the model does not need.

Second, \(\tau_i\) refers to the class of \(x_{i+1}\), an executor that may not
exist, or may not yet be selected, when \(s_i\) emits its continuation. The
condition \(\rho_i(\tau_i)\) is therefore evaluated by the receiving side at
acceptance time, not by the predecessor at emission time. This is the temporal
gap that Proof-of-Continuity addresses for authority~\cite{gallo2026poc},
applied to the envelope: the predecessor fixes the envelope, and any successor
that later appears is admitted only if it falls inside it.

PoR supplies the causal spine. The authority and relational contexts are the two
monotone payloads propagated along that spine.

\section{Dual Monotonicity}

\begin{lemma}[Dual Monotonicity]
Let
\[
\pi=
\langle
(s_0,C_0,\rho_0),
\ldots,
(s_n,C_n,\rho_n)
\rangle
\]
be a chain in which every adjacent transition is relationally valid. Then
\[
C_n\subseteq C_{n-1}\subseteq\cdots\subseteq C_0
\]
and
\[
\rho_n
\preceq_R
\rho_{n-1}
\preceq_R
\cdots
\preceq_R
\rho_0.
\]
In particular,
\[
C_n\subseteq C_0
\]
and
\[
\rho_n\Rightarrow\rho_0.
\]
\end{lemma}

\begin{proof}
Every valid transition satisfies
\[
C_{i+1}\subseteq C_i.
\]
Transitivity of set inclusion gives
\[
C_n\subseteq C_0.
\]

Every valid transition also satisfies
\[
\rho_{i+1}\Rightarrow\rho_i.
\]
Transitivity of logical implication gives
\[
\rho_n\Rightarrow\rho_0.
\]
\end{proof}

The two monotonicities decrease in different mathematical spaces:

\[
C_i
\]
attenuates under set inclusion, while
\[
\rho_i
\]
attenuates under logical implication.

They nevertheless express the same continuity principle:

\begin{quote}
A valid successor may preserve or restrict what its predecessor allowed, but it
may not expand it.
\end{quote}

The proof is short by design. Validity imposes monotonicity at each hop, and
the lemma is its transitive closure, so the result is close to its definition.
The substance of the model lies elsewhere: in the choice of the two ordered
spaces, in their composition with PoR along one lineage, in the reading of the
earlier projection result as a quotient instance, and in the witness
construction of Section~\ref{sec:witness}, which conditionally realizes the
premises under its explicit assumptions.

\section{Relational Closure}

Dual monotonicity implies more than endpoint restriction. Every concrete
transition in the lineage remains inside the relational envelope established at
the origin.

\begin{theorem}[Origin Relational Closure]
Let
\[
\tau_i=
\big(
q_J(x_i),
q_J(x_{i+1})
\big)
\]
be the relational observation of hop \(i\). If every adjacent transition is
relationally valid, then
\[
\tau_i\in\sem{\rho_0}
\]
for every \(i=0,\ldots,n-1\).
\end{theorem}

\begin{proof}
Validity gives
\[
\rho_i(\tau_i)=1,
\]
hence
\[
\tau_i\in\sem{\rho_i}.
\]
By Dual Monotonicity,
\[
\rho_i\Rightarrow\rho_0,
\]
so
\[
\sem{\rho_i}\subseteq\sem{\rho_0}.
\]
Therefore
\[
\tau_i\in\sem{\rho_0}.
\]
\end{proof}

Thus a valid descendant cannot introduce an executor relation that the origin
relational envelope excluded.

\begin{corollary}[No relational expansion]
If
\[
\tau\notin\sem{\rho_0},
\]
then \(\tau\) cannot occur as a valid transition anywhere in the same
relationally continuous lineage.
\end{corollary}

\section{The Product Authorization State}

The two monotone dimensions can be represented as one ordered state.

\begin{definition}[Relational authorization state]
Define
\[
\Gamma_i:=(C_i,\rho_i).
\]
Here and below, the predicate coordinate is its semantic-equivalence class from
Proposition~\ref{prop:lattice}. Thus replacing \(\rho_i\) by a predicate with
the same extension does not change the state.
\end{definition}

\begin{definition}[Product attenuation order]
For states
\[
\Gamma'=(C',\rho')
\qquad\text{and}\qquad
\Gamma=(C,\rho),
\]
define
\[
\Gamma'\preceq\Gamma
\]
iff
\[
C'\subseteq C
\qquad\text{and}\qquad
\rho'\Rightarrow\rho.
\]
\end{definition}

\begin{proposition}[The product order is a partial order]
On
\[
\mathcal P(O\times R)\times
\big(\Pred(\Transitions_J)/{\equiv_{\mathrm{sem}}}\big),
\]
where \(\rho\equiv_{\mathrm{sem}}\sigma\) iff
\(\sem{\rho}=\sem{\sigma}\), the relation \(\preceq\) is a partial order. If
predicates are retained as syntactic formulas instead, the same definition is
only a preorder.
\end{proposition}

\begin{proof}
Reflexivity and transitivity follow componentwise from set inclusion and
logical implication. If both \((C',\rho')\preceq(C,\rho)\) and
\((C,\rho)\preceq(C',\rho')\), then \(C'=C\) and
\(\sem{\rho'}=\sem{\rho}\); hence the two states are equal in the displayed
quotient. Without quotienting formulas by semantic equivalence, distinct but
equivalent formulas witness the failure of antisymmetry.
\end{proof}

Every valid transition therefore satisfies
\[
\Gamma_{i+1}\preceq\Gamma_i.
\]

Hence:
\[
\Gamma_n
\preceq
\Gamma_{n-1}
\preceq
\cdots
\preceq
\Gamma_0.
\]

\begin{lemma}[Combined Closure]
For every valid descendant state \(\Gamma_i=(C_i,\rho_i)\),
\[
C_i\subseteq C_0
\]
and
\[
\sem{\rho_i}\subseteq\sem{\rho_0}.
\]
\end{lemma}

\begin{proof}
Immediate from Dual Monotonicity and the semantic equivalence
\[
\rho_i\Rightarrow\rho_0
\Longleftrightarrow
\sem{\rho_i}\subseteq\sem{\rho_0}.
\]
\end{proof}

The origin therefore bounds both:

\[
\boxed{\text{which privileges may continue}}
\]
and
\[
\boxed{\text{through which executor relationships they may continue}.}
\]

\section{Causal History Grows While Authorization Freedom Shrinks}

The lineage itself is not one of the shrinking objects. Let
\[
\ell_i
:=
\langle s_0,\ldots,s_i\rangle
\]
be the causal prefix ending at step \(s_i\).

Define prefix order
\[
\ell_i\preceq_L\ell_{i+1}
\]
when \(\ell_i\) is a prefix of \(\ell_{i+1}\). Then:
\[
\ell_0
\preceq_L
\ell_1
\preceq_L
\cdots
\preceq_L
\ell_n.
\]

Thus the three dimensions move together but not identically:
\[
\begin{array}{ccccccc}
s_0 &\rightarrow& s_1&\rightarrow&s_2&\rightarrow&\cdots s_n\\[3pt]
C_0&\supseteq&C_1&\supseteq&C_2&\supseteq&\cdots C_n\\[3pt]
\rho_0&\Leftarrow&\rho_1&\Leftarrow&\rho_2&\Leftarrow&\cdots\rho_n\\[3pt]
\ell_0&\preceq_L&\ell_1&\preceq_L&\ell_2&\preceq_L&\cdots\ell_n
\end{array}
\]
where
\[
\rho_{i+1}\Rightarrow\rho_i.
\]

The precise statement is therefore not that the lineage \(\Lineages\) itself
decreases. Rather, causal history grows by prefix extension while the
authorization contexts \emph{carried along that history} attenuate in two
dimensions.

This yields the structural principle:
\[
\boxed{
\text{causal evidence accumulates while authorization freedom attenuates.}
}
\]

\section{The Confused Deputy Under Dual Continuity}

\subsection{Authority mixing}

Suppose an executor \(D\) processes a request whose origin context is \(C_0\)
while simultaneously holding an independent authority source \(C_D\). Let
\[
(o,r)\notin C_0
\qquad\text{but}\qquad
(o,r)\in C_D.
\]

A lineage-invariant possession model may authorize the action by selecting
\(C_D\), even though the request itself did not carry that authority. This is
the authority/causality mismatch that produces the confused deputy in the
multi-step setting~\cite{gallo2026poc}.

\subsection{Valid exercise}

The authorization rule of~\cite{gallo2026poc} links the exercise of a privilege
to the context carried at the step that exercises it. We restate it for the
relational model.

\begin{definition}[Valid exercise]\label{def:exercise}
Let
\[
\pi=
\langle
(s_0,C_0,\rho_0),
\ldots,
(s_n,C_n,\rho_n)
\rangle
\]
be a chain. A step \(s_i\) of \(\pi\) \emph{validly exercises} \((o,r)\) iff
\[
(o,r)\in C_i
\]
and either \(i=0\), or the transition from \(s_{i-1}\) to \(s_i\) is
relationally valid; in the latter case \(\rho_{i-1}(\tau_{i-1})=1\) in
particular.
\end{definition}

The relational condition is indexed by the hop that enters \(s_i\), because an
envelope governs the transition that continues the step carrying it.

\begin{proposition}[No in-lineage exercise beyond the origin]
\label{prop:inlineage}
If every transition of \(\pi\) up to \(s_i\) is relationally valid and \(s_i\)
validly exercises \((o,r)\), then
\[
(o,r)\in C_0.
\]
\end{proposition}

\begin{proof}
By Definition~\ref{def:exercise}, \((o,r)\in C_i\). The prefix
\(\langle(s_0,C_0,\rho_0),\ldots,(s_i,C_i,\rho_i)\rangle\) is itself a chain
whose adjacent transitions are all relationally valid, so Combined Closure
applies to it and gives \(C_i\subseteq C_0\). Hence \((o,r)\in C_0\).
\end{proof}

Independent authority held by an executor does not modify \(C_i\), so it cannot
be imported into the request lineage without violating the authority
attenuation condition. Proposition~\ref{prop:inlineage}, however, only
constrains steps \emph{of} the lineage. It does not yet exclude a deputy that
performs the action \emph{outside} the lineage, by opening a lineage of its own
rooted in its own privileges on the target the request designated. This is
Hardy's compiler, which writes the billing file with its own authority as if it
were housekeeping~\cite{hardy1988confused}. The previous paper identifies this
gap explicitly as lineage origination and leaves it to future
work~\cite{gallo2026poc}. Closing it requires further premises, which we now state
explicitly and separately.

\subsection{Exclusion of the confused deputy}

We first make the objects used by the result explicit. Let \(\mathsf{Req}\) be
the set of requests and \(\mathsf{Eff}\) the set of effects that occur. Every
request \(req\) has a lineage \(\ell_{req}\) and an origin context
\(C_0(req)\). Let
\[
\operatorname{step}:\mathsf{Eff}\rightarrow S,
\qquad
\operatorname{realize}:\mathsf{Eff}\rightarrow O\times R,
\]
associate an effect with the step that performs it and the privilege it
realizes. Let
\[
\operatorname{Lineage}:\mathsf{Eff}\rightarrow\Lineages
\]
be the lineage to which the enforcement system attributes the effect, and let
\[
\operatorname{CausedBy}\subseteq\mathsf{Eff}\times\mathsf{Req}
\]
be the semantic causal relation between effects and requests. Finally, write
\(\operatorname{IndependentAuthority}(D,(o,r))\) when executor \(D\) holds
authority for \((o,r)\) from a source other than the request context. An effect
\(e\) is performed by \(D\) when
\(\exec(\operatorname{step}(e))=D\).

The distinction between \(\operatorname{CausedBy}\) and
\(\operatorname{Lineage}\) is intentional. The former is the causal fact that
the security statement concerns; the latter is the attribution made available
to enforcement.

\begin{definition}[Confused deputy, lineage form]\label{def:cd}
A confused deputy occurs iff there exist
\(req\in\mathsf{Req}\), \(e\in\mathsf{Eff}\), \(D\in X\), and
\((o,r)\in O\times R\) such that
\[
\begin{aligned}
&\operatorname{CausedBy}(e,req),
&&\operatorname{realize}(e)=(o,r),\\
&\exec(\operatorname{step}(e))=D,
&&(o,r)\notin C_0(req),\\
&\operatorname{IndependentAuthority}(D,(o,r)).
\end{aligned}
\]
\end{definition}

The previous paper states condition~1 as \((o,r)\notin Priv(U)\) for the
originating principal \(U\). Since \(C_0\subseteq Priv(U)\), the condition
\((o,r)\notin Priv(U)\) implies \((o,r)\notin C_0\). Every confused deputy in the
sense of~\cite{gallo2026poc} is therefore one in the sense of
Definition~\ref{def:cd}, and excluding the lineage form also excludes the
original one. The lineage form is the stronger target: it also covers
privileges that the originating principal holds but deliberately left out of
\(C_0\).

\begin{definition}[Accepted effect]\label{def:accepted}
An effect \(e\) is \emph{accepted} in lineage \(\ell\) for \((o,r)\), written
\(\operatorname{Accepted}(e,\ell,(o,r))\), iff
\(\operatorname{realize}(e)=(o,r)\),
\(\operatorname{step}(e)=s_k\) for some step \(s_k\) of \(\ell\), every
transition of \(\ell\) up to \(s_k\) has been verified relationally valid, and
\(s_k\) validly exercises \((o,r)\).
\end{definition}

Here \(\operatorname{realize}(e)\) denotes the operation and actual resource on
which the effect occurs after name resolution, aliasing, redirection and other
designation processing. Checking only an input string, pathname or logical
designation is insufficient if it may resolve to a different resource. Complete
mediation requires the verifier to authorize the realized pair \((o,r)\).

\begin{assumption}[Complete mediation]\label{ass:mediation}
Every occurring effect is accepted under its attributed lineage:
\[
\forall e\in\mathsf{Eff},\qquad
\operatorname{Accepted}\bigl(
e,\operatorname{Lineage}(e),\operatorname{realize}(e)
\bigr).
\]
No path to the affected resource bypasses the verifier.
\end{assumption}

\begin{assumption}[Causal attribution]\label{ass:attribution}
Every effect causally triggered by a request is presented as a step of that
request's lineage:
\[
\forall e\in\mathsf{Eff}\;\forall req\in\mathsf{Req},\qquad
\operatorname{CausedBy}(e,req)
\Longrightarrow
\operatorname{Lineage}(e)=\ell_{req}.
\]
\end{assumption}

A third property is related but distinct.

\begin{assumption}[Origination control]\label{ass:origination}
A new lineage is opened only at an origin gate independent of every executor,
and executors hold no authority of their own outside lineages opened there.
\end{assumption}

Origination control and causal attribution are not equivalent. Origination
control restricts \emph{who} may open a lineage, but the gate may still be
invoked because of a request, and would then re-root a request-caused effect in
a new lineage: origination control holds while attribution fails. Conversely,
attribution can hold while executors open genuinely independent lineages of
their own. The theorem below uses complete mediation and causal attribution;
origination control is one ingredient of how an architecture realizes them.

\begin{theorem}[No request-caused effect beyond the origin]\label{thm:causal-bound}
Under Assumptions~\ref{ass:mediation} and~\ref{ass:attribution}, with the
authority-attenuation clause \(C_{i+1}\subseteq C_i\) built into relational
validity, for every
\(req\in\mathsf{Req}\), \(e\in\mathsf{Eff}\), and \((o,r)\in O\times R\),
\[
\operatorname{CausedBy}(e,req)
\land
\operatorname{realize}(e)=(o,r)
\Longrightarrow
(o,r)\in C_0(req).
\]
\end{theorem}

\begin{proof}
Suppose an effect \(e\) caused by \(req\) realizes \((o,r)\). By
Assumption~\ref{ass:attribution}, \(e\) is presented as a step \(s_k\) of
\(\ell_{req}\). By Assumption~\ref{ass:mediation},
\(\operatorname{Accepted}(e,\ell_{req},(o,r))\) holds, so every transition of
\(\ell_{req}\) up to \(s_k\) is relationally valid and \(s_k\) validly
exercises \((o,r)\). Applying Proposition~\ref{prop:inlineage} to the
verified prefix of \(\ell_{req}\) ending at \(s_k\) gives \((o,r)\in C_0\).
Here \(C_0=C_0(req)\), the origin context of \(\ell_{req}\).
\end{proof}

\begin{corollary}[Exclusion of the confused deputy]\label{thm:cd}
Under Assumptions~\ref{ass:mediation} and~\ref{ass:attribution}, the
confused-deputy condition of Definition~\ref{def:cd} is unsatisfiable. A
fortiori, the condition of~\cite{gallo2026poc} is unsatisfiable as well.
\end{corollary}

\begin{proof}
Theorem~\ref{thm:causal-bound} contradicts
\((o,r)\notin C_0(req)\). The independent-authority conjunct is not needed;
the theorem proves the stronger statement that \emph{no} request-caused effect,
whether or not its executor independently holds the privilege, can exercise
authority outside the request's origin context.
\end{proof}

The contradiction is direct, and the enforcement properties that the proof
does not derive sit in two named assumptions. The monotone-authority premise
enters through relational validity and Proposition~\ref{prop:inlineage}; the
relational-envelope premise is not needed for the authority conclusion.
The result does not say that a deputy cannot contain a logic error or attempt an
operation outside \(C_0(req)\). It says that, under complete mediation and
causal attribution, such an error cannot produce an \emph{accepted}
request-caused effect exercising authority outside that context.

\paragraph{What the construction realizes, and what it does not.}
The boundary-witnessed construction of Section~\ref{sec:witness} can discharge
complete mediation for resources whose every access is forced through its
boundaries; exclusive resource mediation is a deployment property, not a
consequence of the checkpoint syntax. Under its Witness Assumptions, the
construction discharges causal attribution for \emph{explicit mediated
handoffs}: a checkpoint is issued only for a step observed consuming the
continuation of a verified predecessor.
For resources subject to that exclusive mediation, origination control also
prevents an executor from opening an independently privileged lineage outside
the gate.

It does not realize attribution for causality that flows through state. A
request may modify an executor's memory, a queue, a database, a cache, a file
or shared memory; the request may then complete, and a later timer, asynchronous
callback or background worker may read that state and obtain a new lineage at
the origin gate. The resulting effect depends causally on the request without
consuming its continuation, and no handoff observation will attribute it.
Designation that flows as data is the typical case: a deputy performs
housekeeping on a target it read from the request. Indirect signaling, covert
channels and side channels can likewise influence control or data flow without
appearing as an explicit mediated handoff. Information-flow tracking or a
stronger execution model can address modeled overt channels, but ordinary
tracking does not by itself establish attribution for every covert or side
channel. Otherwise a further causal-attribution assumption is required, for
example that executors retain no state across requests. The attribution itself
therefore belongs to the trusted computing base of the enforcement
architecture. Corollary~\ref{thm:cd} follows deductively from its two premises;
the construction discharges attribution for mediated handoffs, and not beyond.

\paragraph{Exclusive mediated ingress.}
A strongly confined deployment may impose the sufficient operational condition
that every modeled security-relevant external input capable of influencing a
request-dependent effect enters as a boundary-mediated continuation associated
with a lineage. Combined with complete mediation of effects and
lineage-preserving treatment of every modeled internal channel, this gives the
verifier an ingress and egress point for all such influence. This is a
deployment condition, not a premise of the core closure results, and it does not
establish attribution for unmodeled covert or side channels.

The relational component is not needed for this result. Its contribution is
orthogonal: it closes the class of admissible executor relationships at the
same time.

\begin{corollary}[Dual closure against authority and relationship import]
In a relationally continuous lineage:
\begin{enumerate}
    \item no step validly exercises, and no state carries, a privilege absent
    from \(C_0\); and
    \item no hop has a transition profile outside \(\sem{\rho_0}\), and no
    envelope carried along the lineage admits one.
\end{enumerate}
\end{corollary}

This is the strengthened closure result. Authority cannot be imported from an
unrelated source, and relational admissibility cannot be widened in order to
make an otherwise inadmissible continuation appear valid.

\section{Authorization Without Successor-Side Proof of Possession}

We now separate the semantic authorization rule from the mechanisms used to
produce evidence for that rule.

\subsection{Authentication and authorization are separate layers}

Fixing the steps also fixes their actual executors, since
\(x_i=\exec(s_i)\). To isolate the executor-dependent input, define the
following counterfactual decision template for fixed
\(s_i,C_i,\rho_i,s_{i+1},C_{i+1},\rho_{i+1}\):
\(V:X\times X\rightarrow\{0,1\}\) by
\[
V(x,x')
:=
\PoR(s_i,s_{i+1})
\land
C_{i+1}\subseteq C_i
\land
\rho_i\big(q_J(x),q_J(x')\big)
\land
(\rho_{i+1}\Rightarrow\rho_i),
\]
so that
\[
\ValidR(s_i,C_i,\rho_i,s_{i+1},C_{i+1},\rho_{i+1})
=
V(x_i,x_{i+1}).
\]
\(\PoR\) is a condition on the fixed steps, not on the two arguments of this
template. Thus \(V\) is used only to expose factorization of the relational
input; it does not assert that changing an executor while holding a step fixed
produces another execution admitted by the model.

\begin{proposition}[Relational executor dependence factors through the profile]
\label{prop:layers}
Conditional on fixed execution occurrences and authorization state, the
executor-dependent relational component of the decision factors through
\(q_J\times q_J\): the function \(V\) is invariant under \(\equiv_J\) in each
argument, and therefore factors uniquely as
\[
V=\bar V\circ(q_J\times q_J),
\qquad
\bar V:\Transitions_J\rightarrow\{0,1\}.
\]
\end{proposition}

\begin{proof}
The executors occur in \(V\) only through
\((q_J(x),q_J(x'))\). If \(x\equiv_J y\) and \(x'\equiv_J y'\), then
\(q_J(x)=q_J(y)\) and \(q_J(x')=q_J(y')\), so \(V(x,x')=V(y,y')\). Apply
Proposition~\ref{prop:factorization} to \(Y=X\times X\) with the product
equivalence \(\equiv_J\times\equiv_J\). Its quotient is in canonical bijection
with \(\Profiles_J\times\Profiles_J=\Transitions_J\) via \(q_J\times q_J\), which
is surjective onto \(\Transitions_J\).
\end{proof}

The proposition is deliberately modest. It does not say that the executor is
irrelevant to the decision: the occurrences determine their executors, and
\(\PoR\) depends on the occurrences. By itself it proves invariance under
executor replacement once the profiles and the fixed occurrence-level data are
given; it does not quantify over authentication mechanisms. We now state the
corresponding, and still deliberately limited, mechanism claim explicitly.

For the fixed occurrence and authorization data used above, define
\[
D(u,v):=
\PoR(s_i,s_{i+1})
\land C_{i+1}\subseteq C_i
\land \rho_i(u,v)
\land(\rho_{i+1}\Rightarrow\rho_i),
\qquad (u,v)\in\Transitions_J.
\]

\begin{definition}[Sound profile-evidence mechanism]\label{def:profile-evidence}
A profile-evidence mechanism \(M\) consists of an evidence space
\(\mathcal E_M\) and a partial extraction function
\[
\alpha_M:X\times\mathcal E_M\rightharpoonup\Profiles_J.
\]
It is sound when
\[
\alpha_M(x,e)=u\Longrightarrow u=q_J(x)
\]
for every input on which it is defined. This definition concerns only evidence
for the executor profile; evidence for \(\PoR\) and for external relational
atoms remains subject to its own soundness assumptions.
\end{definition}

\begin{proposition}[Profile-preserving evidence-mechanism invariance]
\label{prop:mechanism-invariance}
Let \(M\) and \(N\) be profile-evidence mechanisms. If
\[
\alpha_M(x,e)=\alpha_N(y,f)=u
\quad\text{and}\quad
\alpha_M(x',e')=\alpha_N(y',f')=v,
\]
then both mechanisms induce the same authorization verdict \(D(u,v)\). In
particular, two sound mechanisms defined for the same concrete executors
\(x_i,x_{i+1}\) induce the same verdict.
\end{proposition}

\begin{proof}
The definition of \(D\) has no argument for a mechanism or for its evidence;
its executor-dependent input is exactly \((u,v)\). Equality of those inputs
therefore gives equality of verdicts. For sound mechanisms applied to the same
executors, Definition~\ref{def:profile-evidence} gives
\(u=q_J(x_i)\) and \(v=q_J(x_{i+1})\) for each mechanism.
\end{proof}

This proposition is an output-preservation result. It does not claim that all
mechanisms can produce evidence, that they have equal security properties, or
that the evidence mechanism for \(\PoR\) may be replaced without re-establishing
its soundness. It states precisely the separation used here: once a sound
mechanism establishes the same profile facts, its internal use or non-use of
PoP is not inspected by \(D\). Possession, where it is used, therefore
contributes evidence to the layer that establishes membership in a relational
class. Holding the same causal and authorization state fixed, the accepted
transition is the same; consequently, the closure results downstream of that
transition are unchanged. This does not transfer the security properties of one
evidence mechanism to another: each must separately satisfy its soundness
obligation. Authentication establishes the executor profiles, PoR establishes
causal continuation, \(C_i\) bounds the privileges that may continue, and
\(\rho_i\) bounds the profile transitions through which they may continue.

Membership is established at admission. The profile is fixed when the boundary
accepts the step and records it in the checkpoint of
Section~\ref{sec:witness}. ``Once authenticated'' therefore means for this
admission, not forever; how class membership changes over time, and how
revocation affects it, remain among the limitations of
Section~\ref{sec:limitations}.

Section~\ref{sec:picx} describes a realization that takes this layered route:
the successor authenticates with possession, and the decision is taken over
what that authentication established.

\subsection{Possession and relationship are different predicates}

A possession predicate has the general form
\[
\operatorname{Possess}(x,k),
\]
meaning that executor \(x\) controls artifact \(k\).

A Proof-of-Possession establishes some proposition of that form, or a closely
related holder-binding proposition:
\[
\PoP(x,k).
\]
It does not by itself establish \(q_J(x)=u\). Inferring profile membership
requires a trusted binding from \(k\) or its credential to the executor and to
the asserted class attributes. Conversely, a sound profile-evidence mechanism
may identify and classify the executor by boundary-local observation without
requiring successor-side PoP.

A Proof of Relationship establishes something categorically different:
\[
\PoR(s_i,s_{i+1}),
\]
namely that one execution occurrence validly continues another.

The two propositions may be connected operationally, but neither is logically
identical to the other.

\subsection{Evidence is not semantics}

A concrete implementation may use a signature, key, credential, attestation,
trusted execution environment, authenticated channel, capability handoff, or
another mechanism to provide evidence for one of the following facts:

\[
q_J(x)=u,
\]
\[
\rho(\tau)=1,
\]
or
\[
\PoR(s_i,s_{i+1}).
\]

PoP may be one such evidence mechanism. For example, possession of a
predecessor-issued continuation artifact may witness a single causal handoff.
But the authorization semantics consumes the proved relationship and the
monotone state; it need not consume possession itself.

\begin{definition}[Possession-independent relational authorization]
\label{def:possindep}
A relational authorization rule is possession-independent when its logical
premises contain no predicate of the form
\[
\operatorname{Possess}(x,k)
\]
or
\[
\PoP(x,k).
\]
\end{definition}

The requirement that PoR and relational-profile facts be soundly verifiable is
not part of the definition; it appears as a hypothesis of
Theorem~\ref{thm:possession-independence}.

Proposition~\ref{prop:layers} places successor-side possession, where it is
used, in the authentication layer. The construction below shows that even that
layer can be discharged without it. This matters for the non-necessity claim:
the rule would be empty if every sound construction of PoR required the
successor to prove possession of something: a non-necessity result has content
only if its premises are satisfiable without the excluded predicate. We
therefore give a conditional construction in which the successor proves
possession of nothing, and state rather than conceal the environmental
assumptions on which it depends.

\subsection{A construction without successor-side possession}
\label{sec:witness}

Let \(h\) be a collision-resistant commitment to the state of an execution
step, and \(H\) a collision-resistant hash. Let each step \(s_i\) be accepted at
an enforcement boundary \(B_i\): the verifier that evaluates \(\ValidR\) for the
hop entering \(s_i\), and that holds a signing key. Write
\(\sigma_B(\operatorname{enc}(m))\) for a signature by \(B\) on the canonical
encoding of a structured payload \(m\), and let \(\lambda\) be a lineage
identifier fixed at the origin. All signed and hashed tuples use the same
injective canonical encoding \(\operatorname{enc}\).

A checkpoint is the tuple
\[
\kappa_i=
\bigl(B_i,m_i,\sigma_{B_i}(\operatorname{enc}(m_i))\bigr).
\]
It is identified by its signer and signed payload, not by the bytes of its
signature. Its identifier is
\[
\operatorname{cid}(\kappa_i)
:=
H\big(\operatorname{enc}(B_i,m_i)\big).
\]
Signature schemes may be malleable, admitting several valid signatures on one
payload; all of them have the same identifier, so malleability affects neither
predecessor binding nor the consumption register.

\begin{definition}[Boundary-witnessed checkpoint]
The checkpoint of step \(s_{i+1}\) is
\[
\kappa_{i+1}
:=
\bigl(
B_{i+1},m_{i+1},
\sigma_{B_{i+1}}(\operatorname{enc}(m_{i+1}))
\bigr),
\]
\[
m_{i+1}
:=
\big(
\operatorname{cid}(\kappa_i),\,h(s_{i+1}),\,\lambda,\,C_{i+1},\,\rho_{i+1},\,
q_J(x_{i+1}),\,\nu_{i+1}
\big),
\]
where \(\nu_{i+1}\) is a nonce generated by \(B_{i+1}\). The boundary
\(B_{i+1}\) issues \(\kappa_{i+1}\) only after:
\begin{enumerate}
    \item verifying the predecessor checkpoint \(\kappa_i\) and that it carries
    the same lineage identifier \(\lambda\);
    \item checking that \(\operatorname{cid}(\kappa_i)\) has not already been
    consumed as the predecessor of an accepted step;
    \item observing that \(s_{i+1}\) consumes the continuation emitted by the
    step committed in \(\kappa_i\), written
    \(\operatorname{ObservedHandoff}(s_i,s_{i+1})\);
    \item identifying the executor \(x_{i+1}\) and determining its profile
    \(q_J(x_{i+1})\); and
    \item checking \(C_{i+1}\subseteq C_i\), obtaining
    \(E_{i+1}\vdash\rho_i(\tau_i)\), and checking
    \(\rho_{i+1}\Rightarrow\rho_i\), where \(E_{i+1}\) is the evidence
    available to the receiving boundary, and where \(C_i\), \(\rho_i\) and
    \(q_J(x_i)\) are read from \(\kappa_i\) and
    \(\tau_i=(q_J(x_i),q_J(x_{i+1}))\). Refinement is checked by the syntactic
    criterion of Section~\ref{sec:refinement}, which is sound.
\end{enumerate}
The origin checkpoint uses the same payload with a distinguished origin marker
\(\bot\) in place of \(\operatorname{cid}(\kappa_{-1})\):
\[
m_0=(\bot,h(s_0),\lambda,C_0,\rho_0,q_J(x_0),\nu_0),
\qquad
\kappa_0=\bigl(B_0,m_0,\sigma_{B_0}(\operatorname{enc}(m_0))\bigr).
\]
Authority bootstrap and the decision to issue \(\kappa_0\) are out of scope as
in~\cite{gallo2026poc}.
\end{definition}

The nonce plays no challenge-response role: no successor answers it. It
distinguishes checkpoints issued by the same boundary for otherwise identical
steps. Since the identifier also covers the issuing boundary \(B_i\), entries in
the consumption register of the anti-replay assumption below are unambiguous
across boundaries as well.

The construction yields a predicate
\[
\PoR_\kappa(s_i,s_{i+1}),
\]
which holds iff \(\kappa_{i+1}\) verifies, its payload's first component is
\(\operatorname{cid}(\kappa_i)\) for a verified \(\kappa_i\), and both carry
\(\lambda\). This
constructed predicate is distinct from the semantic \(\PoR(s_i,s_{i+1})\), which
holds iff \(s_{i+1}\) is a valid causal continuation of \(s_i\). What must be
shown is that \(\PoR_\kappa\) implies \(\PoR\).

Binding \(\operatorname{cid}(\kappa_i)\) rather than only the commitment
\(h(s_i)\) ties each checkpoint to the exact predecessor payload, and with it to the context
\(C_i\), the envelope \(\rho_i\), the profile \(q_J(x_i)\) and the lineage
\(\lambda\) that the checks in item~5 read. Except with negligible probability,
a checkpoint cannot be spliced onto a different predecessor that happens to
share a step commitment.

We work in the computational model. Collision resistance does not make \(H\)
injective: collisions exist, but no efficient adversary finds one except with
probability negligible in the security parameter \(\eta\). Statements about
uniqueness below therefore hold except with negligible probability, not
absolutely.

The construction rests on six separate assumptions.

\begin{assumption}[Witness assumptions]\label{ass:witness}
\leavevmode
\begin{enumerate}
    \item \emph{Unforgeability and binding.} Except with probability negligible
    in \(\eta\), no efficient adversary forges a boundary signature or finds a
    collision in \(h\) or \(H\). Verifiers possess authentic verification keys
    for the boundaries they trust.
    \item \emph{Protocol compliance.} Boundaries follow the protocol: a
    boundary issues a checkpoint only after performing checks 1--5 correctly,
    and records each consumed predecessor.
    \item \emph{Handoff observation and adequacy.} A boundary observes, by its
    own means, which continuation a step consumes. Observed consumption is the
    operational criterion of causal continuation:
    \[
    \operatorname{ObservedHandoff}(s_i,s_{i+1})
    \Longrightarrow
    \PoR(s_i,s_{i+1}).
    \]
    \item \emph{Non-PoP profile observation.} A boundary identifies the
    executor it admits, and soundly determines its profile \(q_J\), from
    boundary-local observation or independently authenticated infrastructure
    evidence, without requiring the successor to demonstrate control of any
    artifact. For example, the boundary may hold a trusted mapping from the
    execution slot, container or virtual machine it controls to the executor,
    and from the executor to its profile:
    \[
    \text{slot}\;\longmapsto\;x_{i+1}\;\longmapsto\;q_J(x_{i+1}).
    \]
    \item \emph{Relational evidence soundness.} For every atom occurring in an
    admitted envelope, the boundary either computes the atom from verified
    profiles or obtains sound, monotone evidence for its truth in the sense of
    Section~\ref{sec:relpred}.
    \item \emph{Anti-replay.} A checkpoint identifier is accepted as the
    predecessor of at most one step, and no boundary reuses a nonce. This requires state: a
    consumption register on which acceptance and marking are atomic, shared by
    the boundaries that may receive a given continuation, or a continuation
    that designates a single verifier with an atomic local register.
\end{enumerate}
``By its own means'' means without relying on credentials presented by the
executor; the paragraph after the proof explains why this matters.
\end{assumption}

\begin{proposition}[Relationship without successor-side possession]
\label{prop:witness}
Under Assumption~\ref{ass:witness}, and except with probability negligible in
\(\eta\), \(\PoR_\kappa(s_i,s_{i+1})\) implies
\(\PoR(s_i,s_{i+1})\). Moreover, if \(\kappa_{i+1}\) verifies, then \(\kappa_i\)
continues into no other accepted step, \(q_J(x_{i+1})\) is the profile of the
executor that realized \(s_{i+1}\), and the conditions \(C_{i+1}\subseteq C_i\),
\(\rho_i(\tau_i)\) and \(\rho_{i+1}\Rightarrow\rho_i\) hold. Verifying
\(\kappa_{i+1}\) involves no predicate \(\operatorname{Possess}(x_{i+1},k)\) for
any artifact \(k\).
\end{proposition}

\begin{proof}
Suppose \(\PoR_\kappa(s_i,s_{i+1})\). By unforgeability, except with
negligible probability, \(B_{i+1}\) issued \(\kappa_{i+1}\). Its payload's first
component is \(\operatorname{cid}(\kappa_i)\) for a verified \(\kappa_i\). A
verified checkpoint on the same payload but with different signature bytes has
the same identifier, and is the same checkpoint for every purpose of the
protocol. A verified checkpoint \(\kappa_i'\) on a different payload or from a
different boundary with \(\operatorname{cid}(\kappa_i')=\operatorname{cid}(\kappa_i)\)
would, since \(\operatorname{enc}\) is injective, be a collision in \(H\): any
execution of the system that produced one would yield an efficient collision
finder, so this happens with negligible probability. Except with
negligible probability, the predecessor binding therefore cannot be redirected
to a different valid checkpoint, and \(\kappa_{i+1}\) is bound to the values in
\(\kappa_i\). By protocol compliance, \(B_{i+1}\) performed checks 1--5 before
issuing. Check~3 together with handoff observation gives
\(\operatorname{ObservedHandoff}(s_i,s_{i+1})\), and handoff adequacy gives
\(\PoR(s_i,s_{i+1})\). Check~2 together with anti-replay gives that
\(\kappa_i\) was consumed by no other accepted step. Check~4 together with
non-PoP profile observation gives that \(q_J(x_{i+1})\) is correct. Check~5
gives \(C_{i+1}\subseteq C_i\) directly and \(\rho_{i+1}\Rightarrow\rho_i\) by
soundness of the syntactic criterion. For \(\rho_i(\tau_i)\), relational
evidence soundness and the lifting rules of Section~\ref{sec:relpred} turn
\(E_{i+1}\vdash\rho_i(\tau_i)\) into \(\rho_i(\tau_i)\). All three
are checked against predecessor values bound in \(\kappa_i\) and successor
values bound in \(\kappa_{i+1}\). Verification
reads \(\kappa_i\), \(\kappa_{i+1}\), the boundaries' public keys, and the
boundary-local or infrastructure evidence permitted by
Assumption~\ref{ass:witness}; nothing in it is held or presented by
\(x_{i+1}\).
\end{proof}

The substantive content of the proposition lies in handoff observation and
adequacy. The
construction converts an observation made by a compliant boundary into
verifiable evidence; it does not create causal knowledge that the boundary did
not have. Protocol compliance makes explicit what the previous paper leaves
implicit when it places compromise of the validator out of
scope~\cite{gallo2026poc}.

\paragraph{Realizing the assumptions without successor-side possession.}
The proposition shows that verification reads nothing from the successor. That
is not enough, because possession could re-enter through the assumptions
themselves. In practice, handoff observation and executor identification are
often realized with a key held by the workload, such as a SPIFFE SVID or an
attestation key presented by the executor. That is a proof of possession by the
successor. If the assumptions were realized that way, the claimed
no-successor-PoP realization would be circular. This is why handoff observation
is required to hold by the
boundary's own means, and why profile determination is required to be non-PoP
profile observation. Two architectural patterns illustrate how this requirement
may be met.
\begin{itemize}
    \item \emph{Launch and measure.} The boundary is the runtime that
    instantiates the executor. It launches \(x_{i+1}\) from an image it has
    measured, and assigns its profile at launch from that measurement and from
    policy. Identity and profile come from the launch record kept by the
    boundary, not from a key held by the workload.
    \item \emph{Co-located observation.} A runtime co-located with the
    executor, such as a hypervisor, a sandbox supervisor, or a proxy through
    which all input reaches the executor, delivers the continuation into the
    executor and observes its consumption directly. It knows which executor it
    hosts because it hosts it.
\end{itemize}
In both, the executor presents nothing: the boundary knows what it observes
because it controls the execution environment. Deployments that instead rely on
workload-held keys use successor-side PoP. The model still applies to them, but
the no-successor-PoP conclusion of
Theorem~\ref{thm:possession-independence} applies only when the Witness
Assumptions are in fact discharged by the deployment.

Possession does not disappear from the system: the boundaries control signing
keys. In this construction it moves from the successor to the witnesses, from
the presenter of authority to the root of trust that attests the causal
handoff. The successor is witnessed; it presents nothing. Verification of a
witness signature establishes the source and integrity of the attestation under
the trusted key; it does not by itself make the attested handoff or profile
true. Those conclusions use handoff adequacy and profile soundness. The
authorization rule then evaluates PoR, \(C_i\), \(\rho_i\), and the established
profile facts.

The same boundary also signs the profile facts \(q_J(x_{i+1})\), in the same
checkpoint and under the same key. Relational authorization therefore requires
no additional checkpoint-signing authority, provided the boundary that PoR
already trusts is also trusted to determine the profile facts. It does not
follow that no further trust is needed. Signing a statement does not make it
true, and the underlying profile evidence, such as the registry of measured
images, the hardware root of the launching platform, or the organizational
authority that assigns classes, may depend on additional trust roots. The
non-PoP profile observation assumption absorbs them.

\begin{theorem}[PoP Independence of the Authorization Rule]
\label{thm:possession-independence}
Assume:
\begin{enumerate}
    \item \(\PoR(s_i,s_{i+1})\) is soundly verifiable;
    \item quotient-profile membership \(q_J(x)\) is soundly verifiable;
    \item relational predicates are soundly evaluable; and
    \item authority and relational contexts satisfy their respective
    monotonicity conditions.
\end{enumerate}
Then relational authorization can be decided by a possession-independent rule in
the sense of Definition~\ref{def:possindep}. Moreover, conditional on
Assumption~\ref{ass:witness}, the checkpoint construction realizes these
premises without requiring any successor executor to prove possession of an
artifact; the cryptographic possession used by that mechanism resides in the
boundary witnesses' signing keys. This second conclusion is relative to the
Witness Assumptions, not an unconditional proof that every required
environmental observation can be realized in every deployment.
\end{theorem}

\begin{proof}
Under the assumptions, a transition is accepted exactly when
\[
\PoR(s_i,s_{i+1}),
\]
\[
C_{i+1}\subseteq C_i,
\]
\[
\rho_i(\tau_i),
\]
and
\[
\rho_{i+1}\Rightarrow\rho_i
\]
all hold.

No possession predicate occurs in this decision rule. This first conclusion is
therefore immediate from the definition of the rule once assumptions 1--4 are
granted. For the conditional realization,
Proposition~\ref{prop:witness} gives, under Assumption~\ref{ass:witness}, a
construction in which \(\PoR_\kappa\) soundly witnesses \(\PoR\) and which
establishes assumptions (1), (2) and (4) at the boundary. Assumption (3) holds
by relational evidence soundness: the boundary computes each atom from the
verified profiles or holds sound evidence for it, and soundness lifts to every
positive predicate. For this conditional construction, it suffices to restrict
envelopes to atoms decidable from the profiles, such as equality and inequality
of profile components. Handoff observation by the boundary's own means and
non-PoP profile observation are hypotheses, so no successor proves possession
of anything within the construction. Therefore possession is not logically
necessary to the authorization semantics, even though a concrete implementation
may choose to use PoP as one mechanism for establishing some of the premises.
\end{proof}

The launch-and-measure and co-located patterns above are architectural examples,
not a formal witness model proving the joint satisfiability of
Assumption~\ref{ass:witness}. Accordingly, the checkpoint argument is a proof
sketch in the computational model whose soundness is conditional on the six
listed assumptions.

Theorem~\ref{thm:possession-independence} concerns admission of a transition by
\(\ValidR\). It yields a statement about an occurring resource effect only when
complete mediation connects effects to accepted steps; claims about effects
caused by a request additionally use causal attribution, exactly as in
Theorem~\ref{thm:causal-bound}.

\subsection{What the result does and does not claim}

The theorem is a non-necessity result, not a claim that PoP has no security
role.

PoP may still be useful for:
\begin{itemize}
    \item proving control of a private key;
    \item holder binding;
    \item authenticating an evidence producer;
    \item proving control of a continuation artifact;
    \item protecting transport or replay-sensitive mechanisms; and
    \item implementing a particular PoR construction.
\end{itemize}

The distinction is semantic:
\[
\boxed{
\text{PoP may provide holder-binding or authentication evidence;}
}
\]
whereas
\[
\boxed{
\begin{array}{c}
\text{PoR establishes causal continuation;}\\
\text{\(C_i\) bounds authority and \(\rho_i\) bounds relational continuation.}
\end{array}
}
\]

Within this model, possession presented by the successor is therefore an
optional evidence or authentication-layer mechanism rather than a necessary
premise of the authorization rule. The conditional witness construction places
cryptographic possession at the boundary rather than at the successor; other
deployments may use PoP at either layer without making it a semantic premise.

\begin{remark}[Mechanism independence]\label{rem:mechanism}
Outside the witness construction of Section~\ref{sec:witness}, which is one
realization among many, no result in this paper refers to cryptography.
\(\PoR\), profile membership and the monotone state enter only as predicates;
signatures, hashes, attestation and proofs of possession are ways of producing
evidence for them. Proposition~\ref{prop:mechanism-invariance} proves invariance
under replacement of a profile-evidence mechanism only when the replacement is
sound and supplies the same profile. It does not prove that arbitrary evidence
mechanisms are interchangeable, nor does it remove the separate soundness
obligation for evidence of \(\PoR\) or relational atoms. PoR and PoC are modeled
as properties of the execution; PoP is one possible property established by an
evidence protocol.
\end{remark}

\section{From Possession to Relation}

A possession-centered authorization rule can be abstracted as
\[
\operatorname{Possess}(x,a)
\land
\operatorname{Grants}(a,o,r)
\Longrightarrow
\operatorname{Authorized}(x,o,r).
\]

The relational model instead evaluates an execution occurrence:
\[
\PoR(s_i,s_{i+1})
\]
together with monotone state evolution
\[
\Gamma_{i+1}\preceq\Gamma_i.
\]

The authorization question therefore changes from:
\begin{quote}
What authority does this actor possess?
\end{quote}
to:
\begin{quote}
Is this occurrence a valid continuation of an already authorized causal
relation, under an authorization state that has not expanded?
\end{quote}

This does not eliminate identity, credentials, or possession mechanisms.
Rather, it relocates them. Identity and PoP may establish facts about the
executor or the evidence; the authorization relation is evaluated over the
causal and relational structure of the execution.

\section{A Realization of the Premises: PIC-X}
\label{sec:picx}

The results of this paper are stated over abstract predicates and do not depend
on any implementation. This section records how one open-source implementation,
PIC-X~\cite{picx}, part of the PIC protocol~\cite{picprotocol}, discharges the
restricted subset of premises identified below. Its purpose is an engineering
illustration, not a proof or a complete witness of the model. The description was checked against the
PIC-X source at tag \texttt{v0.2.2}, which builds on the \texttt{pic-protocol}
crate, version \texttt{0.2.3}.

PIC-X takes the layered route of Proposition~\ref{prop:layers}, not the witness
route of Section~\ref{sec:witness}. The successor authenticates with
possession, and each checkpoint is signed by the realm that settles the
transition.

\paragraph{A difference in terminology.}
PIC-X uses the name ``Proof of Relationship'' for the evidence that
authenticates the workload: an SD-JWT~\cite{sdjwt} issued by an attester that
the realm trusts, which binds the workload's public key through its
confirmation claim. In the terms of this paper, that evidence belongs to the
authentication layer. The causal relationship that this paper calls \(\PoR\) is
realized in PIC-X by the transition, which references the exact predecessor
checkpoint.

\paragraph{Settlement.}
The workload submits a candidate transition and signs it, together with the
candidate continuity and the candidate token, with the key that the attester
bound. The realm then verifies, in order: the attester's presentation and the
workload's three signatures; that the predecessor checkpoint verifies under the
realm's own key; the hash of the exact predecessor bytes, the position
progression, the challenge chain, the lineage identifier and the expiry; the
conformance of the disclosed claims to the execution contract; non-expansion;
revocation; and local policy. Only then does it issue the next checkpoint, signed
by the realm.

\begin{table}[ht]
\centering
\small
\begin{tabular}{@{}p{0.27\linewidth}p{0.67\linewidth}@{}}
\toprule
Model element & PIC-X \texttt{v0.2.2} \\
\midrule
Profile \(q_J(x)\) &
Claims disclosed in the attester-issued SD-JWT presentation, verified at
settlement and recorded with the attester's identity for attribution. \\
Relation \(\sim_j\) &
Not a separate object. Equality of the disclosed value of claim \(j\), from a
trusted attester, is the induced equivalence: the kernel of the claim function. \\
Authentication of membership &
Attester-issued SD-JWT binding the workload key in its confirmation claim; the
workload signs the transition, the candidate continuity and the candidate token
with that key. Key-binding JWTs are not used, and are rejected. \\
\(\PoR(s_i,s_{i+1})\) in this paper &
Only partially realized. The transition carries the SHA-256 hash of the exact
predecessor checkpoint bytes, position \(n+1\), challenge continuity and the
checkpoint's lineage identifier, and the predecessor checkpoint must verify
under the realm's own key. This evidences possession of the predecessor
checkpoint, not an observed handoff: handoff observation and adequacy are not
realized. \\
\(C_i\) and \(C_{i+1}\subseteq C_i\) &
Context of authority: identity context and invariants, attenuated by removal
bitmaps and checked by subset inclusion. \\
\(\rho_i\) and \(\rho_{i+1}\Rightarrow\rho_i\) &
Execution contract: key--value and membership constraints on the successor's
disclosed claims. The predecessor's contract governs admission; additions are
conjunctive only, and no entry can be removed or weakened. Unary and
conjunctive; binary atoms and disjunction are not implemented. \\
Checkpoint signing &
Each checkpoint is a COSE object signed by the realm; accepting a predecessor is
a signature check under the realm's published keys. \\
Freshness &
SD-JWT validity claims, with bounded clock skew. A checkpoint stays advanceable
while the realm key that signed it is published. The protocol crate provides a
revocation hook at settlement, but PIC-X \texttt{v0.2.2} configures none. No
credential status list in this version. \\
Linearity and replay &
Stateless: a checkpoint is not retired when a successor is issued, so fan-out
is supported. There is no single-consumption register. \\
\bottomrule
\end{tabular}
\caption{Correspondence between the model and PIC-X \texttt{v0.2.2}, with
\texttt{pic-protocol} \texttt{0.2.3}.}
\label{tab:picx}
\end{table}

\paragraph{The execution contract as a unary envelope.}
The execution contract is a restricted envelope. Its atoms depend only on the
successor's profile, \(\rho_i(u,v)=\bigwedge_k R_k(v)\), and they are combined by
conjunction only. Addition-only accumulation is exactly the syntactic refinement
\(\rho_{i+1}=\rho_i\land\sigma\) of Section~\ref{sec:refinement}. The contract
is evaluated from the predecessor checkpoint, which matches the design choice
that \(\rho_i\) governs hop \(i\). Atoms that relate predecessor and successor,
such as the separation-of-duty atom of Section~\ref{sec:sod}, are beyond this
version.

\paragraph{Why positive predicates suit selective disclosure.}
The evidence monotonicity of the positive language has a direct reading here. A
holder that withholds a disclosure removes evidence, and removing evidence
cannot produce an admission that full disclosure would deny. Selective
disclosure therefore cannot be used to widen admission. PIC-X behaves this way:
a contract entry that is not disclosed is a rejection. The same property
explains why constraints must be positive attested claims: an SD-JWT cannot
establish the absence of a claim by withholding it.

\paragraph{Partial profiles.}
Selective disclosure makes the verified profile partial: the boundary learns
only the components that were disclosed, while \(q_J(x)\) in the model is total.
The gap is the one the evidence relation \(\vdash\) addresses. The boundary
admits on \(E\vdash\rho_i(\tau_i)\), where \(E\) contains the verified
disclosures. Soundness of \(E\) gives \(\rho_i(\tau_i)\) for the true profile,
and undisclosed components can only prevent admission.

\paragraph{Where the realization departs from the model.}
The following departures affect no theorem, but they bound what PIC-X
\texttt{v0.2.2} realizes.
\begin{itemize}
    \item \emph{Successor-side possession.} PIC-X authenticates the successor
    with a key it holds. This is the layered route of
    Proposition~\ref{prop:layers}: possession sits in the authentication layer
    and does not enter the decision. The conditional no-successor-PoP
    realization of Theorem~\ref{thm:possession-independence} is the construction
    of Section~\ref{sec:witness}, not PIC-X.
    \item \emph{Fan-out.} PIC-X admits sibling branches from one checkpoint,
    while this paper models linear chains. Every root-to-leaf path is a linear
    chain, and every result of this paper applies along it. No branch can import
    authority from another, because each successor is attenuated from its own
    predecessor. The anti-replay assumption of the witness construction does not
    hold in PIC-X, and is needed only for the claim that a checkpoint continues
    into no other accepted step, which PIC-X deliberately does not provide.
    \item \emph{Relationship as possession of the checkpoint.} PIC-X does not
    observe a handoff. The relationship is evidenced by the workload's signature
    over a transition that references the exact predecessor bytes and continues
    its challenge. The checkpoint travels inside the settled token, and its
    challenge is in clear. Any workload that is attested by a trusted attester,
    conforms to the execution contract and obtains the settled token can
    therefore advance the lineage, whether or not the token was handed to it.
    The damage is bounded by the checkpoint's authority and by the contract,
    but the relationship reduces to possession of a continuation artifact. With
    an unknown successor, nothing placed in the checkpoint at emission time can
    close this gap statelessly; observing the handoff at consumption, as in the
    construction of Section~\ref{sec:witness}, is left to a separate component.
    \item \emph{No request binding.} The protocol crate offers a
    request-binding hook, but PIC-X leaves it at its default, which accepts
    every transition. A checkpoint is therefore not tied to the request it
    continues.
    \item \emph{Unconstrained settled token.} The settled token carries no
    audience and no key binding. With an unknown successor this is
    unavoidable, but it makes the token a bearer credential for every attested,
    contract-conforming workload.
    \item \emph{No revocation.} No revocation check is configured; a checkpoint
    is bounded only by the lineage expiry and by the retention of the realm key.
    \item \emph{Attribution.} Complete mediation and causal attribution, the
    premises of Theorem~\ref{thm:causal-bound} and
    Corollary~\ref{thm:cd}, are properties of the deployment around PIC-X,
    such as resource servers that accept only settled tokens, not of PIC-X
    alone. Given the first item, PIC-X does not by itself supply causal
    attribution, including for dependencies carried through persistent state,
    asynchronous work, covert channels or side channels.
\end{itemize}

\paragraph{What the realization does not change.}
No theorem of this paper depends on PIC-X. PIC-X realizes authority
attenuation, a unary conjunctive envelope, and profile soundness resting on the
attesters trusted for each claim; the linear results apply per path. It does
not realize handoff observation and adequacy, anti-replay, request binding, or
revocation, binary relational atoms, disjunction, complete mediation, or
side-channel attribution, so it does not discharge the Witness Assumptions of
Section~\ref{sec:witness} nor, by itself, the premises of
Theorem~\ref{thm:causal-bound}.

\section{Related Work}
\label{sec:related}

\paragraph{Capabilities.}
Capabilities fuse designation and authority at invocation and thereby prevent
the classical confused deputy when authority is exercised through the supplied
capability rather than through ambient authority
\cite{dennis1966programming,miller2003capability}. In the terminology used
here, such an invocation may realize a strong single-hop relationship between
designation and authority.

The present model is compatible with capability implementations. A capability
may be one way to establish a relational predicate, and an attenuated,
context-bound capability may carry both authority and relational information.
The distinction is not where the evidence is encoded but whether the
authorization decision is invariant under changes in causal occurrence.

\paragraph{Macaroons and caveated tokens.}
Macaroons attenuate a bearer credential by appending caveats, and caveats can
confine authority by context, including by which party may exercise
it~\cite{birgisson2014macaroons}. Caveat chains are monotone in the same sense
as \(C_i\) and \(\rho_i\): adding a caveat can only restrict. A macaroon whose
caveats bind the lineage and the executor class would therefore carry both
monotone dimensions inside one artifact. In the terms of this paper it is a
possession-based encoding of the envelope. Its continuity guarantees hold
exactly to the extent that its caveats are lineage-sensitive, which is the
condition identified in~\cite{gallo2026poc}. The present model states that
condition independently of the artifact that carries it.

\paragraph{Attenuated delegation chains.}
Monotone attenuation along a chain is not itself new. Biscuit tokens carry
Datalog checks, and the holder may append blocks that only add checks,
restricting the token's rights~\cite{biscuitspec}. UCAN delegation chains
require each delegation to restate or attenuate the capabilities of its
delegator~\cite{ucanspec}. Two recent individual Internet-Drafts apply the idea
to AI agents. One profiles OAuth JWT access tokens so that each delegation
attenuates scope and constraints, with parent and child linked by a hash
commitment and verifiable offline~\cite{asor2026delegation}. The other binds
attenuated capability grants to encrypted task capsules, attenuating decryption
access together with authorization~\cite{atakora2026sadp}. Hash-linked
predecessor binding, authority attenuation along a chain, monotone accumulation
of restrictions and lineage verification are therefore shared with this line of
work; the witness construction of Section~\ref{sec:witness} uses them and does
not claim them as contributions. These systems attenuate delegated authority
encoded in an artifact. The contribution here is a second ordered dimension,
\(\rho_{i+1}\Rightarrow\rho_i\) over relational transition profiles of
executors, separate from the carrier of authority and composed with authority
attenuation and causal continuity, together with the semantic separation from
successor-side PoP.

\paragraph{Attribute-based access control.}
ABAC authorizes an operation by evaluating attributes of the subject, the
object, the operation and the environment against policy~\cite{hu2014abac}.
The profile map \(q_J\) is an ABAC-style classification of executors, and atoms
over \(\Transitions_J\) are attribute conditions over pairs of executors. What
the model adds is not the classification but its propagation. An ABAC decision
point evaluates the policy attached to the resource it guards, with no
guarantee that a downstream decision point applies a policy no weaker than an
upstream one. Here the envelope \(\rho_i\) is carried by the execution, bound
to its lineage by PoR, and can only be refined along it. Relational closure is
a property of the lineage, not of any single decision point.

\paragraph{Separation of duty.}
The NIST RBAC model standardizes dynamic separation of duty as a constraint on
the roles a user may activate within a session~\cite{ferraiolo2001rbac}, and
history-based formulations make separation of duty depend on which principals
have already acted~\cite{simon1997sod}. The adjacent case of
Section~\ref{sec:sod} is a lineage-scoped variant: the constraint is carried by
the execution rather than attached to a session or an object, and no descendant
can drop it. The non-adjacent case needs memory of earlier executors, which is
the territory of the history-based formulations.

\paragraph{History-based access control.}
History-based access control likewise makes execution history relevant to
authorization~\cite{abadi2003history}. The construction here differs in
emphasis: it treats executor equivalence classes and relational predicates as a
first-class authorization space, then propagates that space monotonically along
a PoR-linked discrete execution.

\paragraph{Proof-carrying and distributed authorization.}
In proof-carrying authorization, a requester presents a proof, in an
authorization logic over credentials, that it is entitled to access, and a
reference monitor checks the proof~\cite{appel1999pca,bauer2005distributed}.
The object of such a proof is the requester's authority, typically rooted in
credentials it holds or can obtain. The object of PoR is different: a
relationship between two execution occurrences. The two are complementary. A
proof-carrying system could serve as the mechanism that verifies the premises of
\(\ValidR\), and the lineage state \((C_i,\rho_i)\) could appear as a premise in
such a logic.

\paragraph{Scope of the quotient construction.}
The quotient construction is also deliberately agnostic to the concrete
meaning of each equivalence relation. One relation may encode trust-domain
membership, another a hardware-attestation class, another a workload class,
and another an organizational relationship. Properties that depend on
occurrences rather than executors, such as an execution-specific continuation
property, do not belong in an equivalence relation on \(X\); they live in
\(\PoR\) or in the carried state. The authorization semantics depends only on the
resulting quotient classes and relational predicates.

\section{Discussion}

\subsection{Why quotient spaces matter}

The quotient construction is not merely notation. It separates concrete
identity from authorization-relevant identity.

Suppose \(x\neq y\) but
\[
x\sim_j y.
\]
Then any policy intentionally defined only over
\[
\Rel_j=X/{\sim_j}
\]
must treat \(x\) and \(y\) equally. If the policy must distinguish them, then
\(\sim_j\) was too coarse and a finer relational view must be introduced.

Multiple quotient spaces therefore provide a disciplined way to state which
executor differences matter to authorization and which do not.

\subsection{Refinement instead of replacement}

A lineage need not carry one fixed relational predicate forever. It may refine
the predicate:
\[
R_1\lor R_2
\;\succeq_R\;
R_1
\;\succeq_R\;
R_1\land R_3.
\]

This is the relational analogue of authority attenuation:
\[
\{\text{read},\text{write}\}
\supseteq
\{\text{read}\}.
\]

The shared principle is non-expansion.

\subsection{Dynamic separation of duty}
\label{sec:sod}

The relational envelope constrains which executors may act along a lineage, not
only which authority they carry. It therefore expresses constraints between
executors that authority attenuation alone cannot state.

Let \(\sim_a\) classify executors by role, and let \(A\subseteq\Rel_a\) be the
approver classes. For a relation \(\sim_j\) and profiles \((u,v)\in\Transitions_J\),
with \(u_j\) the \(j\)-th component of \(u\), define the atom
\[
\mathrm{SoD}_j(u,v)
:\Longleftrightarrow
v_a\notin A
\;\lor\;
u_j\neq v_j.
\]
The atom is a fixed subset of \(\Transitions_J\) and is attested like any other.
It holds trivially when the successor is not an approver; when it is, the
approver must lie in a different \(\sim_j\)-class from its predecessor.

A lineage whose proposal step \(s_i\) is followed by an approval step can refine
its envelope at the proposal step to
\[
\rho_i:=\rho_{i-1}\land\mathrm{SoD}_j.
\]
This is a monotone refinement, checked by the receiving boundary at acceptance.
Because every later envelope implies \(\rho_i\), no descendant can drop the
constraint. The result is dynamic separation of duty, the four-eyes principle,
expressed as a refinement of the relational envelope.

Authority attenuation cannot state this constraint: proposer and approver may
carry the same authority context, and the constraint concerns who exercises it.
A pure bearer capability, without holder or context caveats, is monotone in
authority and silent on the subject that exercises it, so the constraint is not
a property of such a capability. Caveated tokens can express it, as discussed in
Section~\ref{sec:related}. Here it is one refinement of \(\rho\).

The atom relates adjacent executors, and the adjacent case lies within the
formalism. The non-adjacent case does not. A predicate sees only
\(\tau_i=(q_J(x_i),q_J(x_{i+1}))\), and a proposer class recorded in a
checkpoint is a state variable outside that domain. A formal treatment extends
the state to
\[
\Gamma_i=(C_i,\rho_i,m_i),
\]
where \(m_i\) is a write-once memory attested by the boundary, such as the
proposer's class, and lets predicates range over \(M\times\Transitions_J\). The
closure arguments use only extension inclusion and transitivity, so they carry
over to the extended domain unchanged. We leave that development to future
work.

\subsection{The role of causality}

Neither authority attenuation nor relational refinement alone establishes
continuity. The two sequences could be monotone but belong to unrelated
executions.

PoR is what binds the monotone states to one causal history:
\[
s_0\rightarrow s_1\rightarrow\cdots\rightarrow s_n.
\]

Thus the model has three inseparable components:
\begin{enumerate}
    \item causal linkage by PoR;
    \item monotone privilege authority; and
    \item monotone relational admissibility.
\end{enumerate}

\subsection{The lineage axis}

Previous work represented an authority occurrence as
\[
(o,r,\ell)\in O\times R\times\Lineages,
\]
where \(\ell\) identifies the causal lineage. The present model does not replace
that axis.

Instead, it refines the structure carried \emph{along} a lineage. The lineage
prefix
\[
\ell_i
\]
grows, while the attached state
\[
(C_i,\rho_i)
\]
shrinks in the product attenuation order.

It is therefore mathematically more precise to say:
\begin{quote}
continuity is monotone in the authorization state carried along the lineage,
not that the lineage itself is decreasing.
\end{quote}

\section{Limitations and Future Work}
\label{sec:limitations}

The model is intentionally minimal and several questions remain open.

First, the equivalence relations \(\sim_j\) are treated as given. A practical
system must define who is authorized to assert them, how evidence is verified,
how classes change over time, and how revocation affects class membership.

Second, the paper assumes a common relational vocabulary \(J\). Heterogeneous
systems may use different executor classifications at different boundaries. A
general treatment would require sound translations between quotient spaces,
analogous to authority translation between heterogeneous operation/resource
spaces.

Third, the model is linear. Forks, joins, fan-out, fan-in, and DAG-shaped
executions require a partial-order generalization of both authority attenuation
and relational refinement. The anti-replay assumption of the witness
construction, which lets a checkpoint precede at most one step, excludes
fan-out by construction and is consistent with this linear scope. A DAG
generalization must replace it with a rule stating how one checkpoint may be
consumed by several successors.

Fourth, PoR remains largely abstract. The paper proves properties conditional
on sound verification of the causal relationship. The boundary-witnessed
construction of Section~\ref{sec:witness} gives a no-successor-PoP realization
relative to its Witness Assumptions, but does not independently prove their
joint satisfiability. It is only a computational proof sketch and is not
prescribed as the only construction.

Fifth, the exclusion of the confused deputy rests on complete mediation and
causal attribution (Assumptions~\ref{ass:mediation}
and~\ref{ass:attribution}). The model does not decide whether an effect was
caused by a request. The boundary-witnessed construction realizes attribution
for explicit mediated handoffs only; causality that flows through mutable
memory, queues, databases, storage, caches, shared files, timers, asynchronous
callbacks or background workers requires a stronger execution model,
information-flow tracking for the modeled channels, or a further assumption.
Covert channels and side channels are not resolved by the witness construction:
they may influence control or data flow without appearing as a mediated
handoff, and ordinary information-flow tracking need not capture all of them.
The confused-deputy result remains valid conditional on Causal Attribution, but
the construction does not establish that assumption for arbitrary hidden
channels. An error in attribution is a failure of the trusted computing base,
not a violation of the continuity invariant.

The witness construction also separates assumptions that deployments must
realize independently: handoff observation, non-PoP profile observation,
whose evidence may rest on further trust roots, and anti-replay, which requires
shared state among the boundaries that may receive a continuation. Dynamic
separation of duty between non-adjacent steps requires the state extension
sketched in Section~\ref{sec:sod}.

Sixth, the possession-independence result is semantic. A deployment may still
choose PoP for engineering reasons, including authentication, key confirmation,
replay resistance, or evidence transport. The result is only that PoP is not a
necessary logical premise once the required relationship facts are available by
some sound means.

Finally, the model does not claim that every authorization policy should be
relational. It identifies a relational authorization structure for discrete
causal executions where authority continuity matters.

\section{Conclusion}

Proof-of-Continuity established that authority can be propagated through a
causal execution by combining Proof of Relationship with monotone authority
attenuation. This paper develops the relational structure implicit in that
model.

Executors are classified through authorization-relevant equivalence relations:
\[
\sim_j\subseteq X\times X.
\]
Each relation induces a quotient:
\[
\Rel_j=X/{\sim_j},
\]
which identifies physically different executors whenever their differences are
irrelevant to the selected authorization characteristic.

Multiple relational quotients coexist, and their combined profile
\[
q_J(x)
\]
provides the relational coordinates of an executor. Predicates over transitions
between such profiles compose with conjunction and disjunction, and implication
defines an attenuation order:
\[
\rho_{i+1}\Rightarrow\rho_i.
\]

A continuous execution therefore carries two monotone authorization dimensions:
\[
C_{i+1}\subseteq C_i
\]
and
\[
\rho_{i+1}\Rightarrow\rho_i.
\]

At the same time, causal history accumulates:
\[
\ell_i\preceq_L\ell_{i+1}.
\]

The resulting structure is:
\[
\boxed{
\text{history grows while authority and relational freedom shrink.}
}
\]

Dual monotonicity gives combined closure. No valid descendant may introduce a
privilege absent from the origin authority context, and no valid hop may
introduce a relational transition excluded by the origin relational envelope.
Under two explicit assumptions, complete mediation and causal attribution, the
request-bounded-effect theorem excludes any accepted request-caused effect
outside the origin authority context, and the confused-deputy result follows as
a corollary. This does not preclude a buggy deputy from attempting such an
operation; it precludes acceptance of the resulting effect under the stated
premises. The enforcement boundary construction covers attribution only for
explicit handoffs under its Witness Assumptions, and not for arbitrary state,
covert or side channels. The authorization model is at the same time
strengthened with relational closure.

Most importantly, the logical authorization rule contains no necessary
successor-side possession premise. If PoR and relational-profile facts are
soundly verifiable, PoP may supply holder-binding evidence or be used by one
particular witness mechanism, but it is not required as the semantic basis of
transition admission. The checkpoint construction supports a no-successor-PoP
realization only relative to its explicit Witness Assumptions. Claims about
resource effects additionally require complete mediation and causal
attribution.

The distinction can be stated compactly:
\[
\boxed{
\text{possession can evidence control; PoR establishes causal continuity.}
}
\]
Equivalently, in terms of the two layers:
\[
\boxed{
\begin{array}{c}
\text{sound evidence establishes executor profiles;}\\
\text{PoR identifies the causal continuation;}\\
\text{\(C_i\) bounds authority and \(\rho_i\) bounds relational continuation.}
\end{array}
}
\]

This moves authority from a model centered on what an executor possesses to a
model centered on which causal and relational continuations remain valid.

\section*{Acknowledgment of Prior Model}

This paper extends the Proof-of-Continuity and Provenance Identity Continuity
model developed in earlier work~\cite{gallo2026poc}. The new
contribution is the executor quotient construction, the relational predicate
order, dual monotonicity, relational closure, the reading of the earlier
projection result as a quotient instance, and the possession-independence
result together with a conditional construction that satisfies its premises
without successor-side possession when the Witness Assumptions hold.

\appendix
\section{Worked Example}
\label{app:example}

Let \(J=\{\mathrm{org},\mathrm{role},\mathrm{prin}\}\), where
\(\sim_{\mathrm{prin}}\) classifies executors by the principal on whose behalf
they act. Several executor instances may share a principal, so its classes are
not singletons. Profiles are triples
\(u=(u_{\mathrm{org}},u_{\mathrm{role}},u_{\mathrm{prin}})\), and the approver
classes are \(A=\{\mathrm{approver}\}\subseteq\Rel_{\mathrm{role}}\). Define the
atom
\[
\mathrm{SameOrg}(u,v):\Longleftrightarrow u_{\mathrm{org}}=v_{\mathrm{org}},
\]
and take \(\mathrm{SoD}_{\mathrm{prin}}\) from Section~\ref{sec:sod}. Both atoms
are decidable from the profiles.

\paragraph{The lineage.}
A request enters at \(s_0\), realized by \(x_0\) with profile
\((\mathrm{acme},\mathrm{requester},\mathrm{alice})\), with
\[
C_0=\{(\mathrm{read},r),(\mathrm{write},r),(\mathrm{approve},p)\},
\qquad
\rho_0=\mathrm{SameOrg}.
\]
Step \(s_1\) is realized by \(x_1\) with profile
\((\mathrm{acme},\mathrm{approver},\mathrm{bob})\). The hop from \(s_0\) is
admitted, since \(\rho_0\) holds on \((q_J(x_0),q_J(x_1))\). At \(s_1\), bob
proposes a change and refines the envelope:
\[
C_1=C_0,
\qquad
\rho_1=\mathrm{SameOrg}\land\mathrm{SoD}_{\mathrm{prin}}.
\]
The refinement \(\rho_1\Rightarrow\rho_0\) holds syntactically, by conjunction.
Step \(s_2\) is an approval by executor \(x_2\). It presents evidence, for
instance an SD-JWT from an issuer trusted for organization, role and principal,
disclosing \(\mathrm{org}=\mathrm{acme}\), \(\mathrm{role}=\mathrm{approver}\)
and \(\mathrm{prin}=\mathrm{carol}\). The boundary verifies it and fixes
\(q_J(x_2)=(\mathrm{acme},\mathrm{approver},\mathrm{carol})\).

\paragraph{Acceptance.}
With \(\tau_1=(q_J(x_1),q_J(x_2))\), the boundary checks \(\rho_1(\tau_1)\). The
organizations agree. The successor is an approver, and its principal differs from
the predecessor's, since \(\mathrm{bob}\neq\mathrm{carol}\). Both conjuncts
hold. The boundary also checks \(C_2=\{(\mathrm{approve},p)\}\subseteq C_1\)
and \(\rho_2=\rho_1\), a trivial refinement. The hop is accepted. By
Proposition~\ref{prop:layers}, any executor with the profile of \(x_2\), for
instance another process acting for carol, would receive the same verdict,
however it had authenticated.

\paragraph{Rejections.}
Four variants fail, each for a different reason.
\begin{enumerate}
    \item Another executor instance acting for bob approves:
    \(\tau=((\mathrm{acme},\mathrm{approver},\mathrm{bob}),
    (\mathrm{acme},\mathrm{approver},\mathrm{bob}))\). The successor is an
    approver and the principals coincide, so \(\mathrm{SoD}_{\mathrm{prin}}\)
    fails. Separation is between principals, not instances: a second process of
    the same principal is not a second pair of eyes.
    \item An approver from another organization: \(\mathrm{SameOrg}\) fails.
    \item An approver that withholds the organization disclosure: the evidence
    does not support \(\mathrm{SameOrg}\), so \(E\nvdash\rho_1(\tau_1)\) and the
    hop is rejected. By evidence monotonicity, no withholding can turn a
    rejection into an admission.
    \item An approver whose transition carries
    \((\mathrm{delete},r)\notin C_1\): non-expansion fails.
\end{enumerate}
In none of these cases does holding a valid credential change the outcome: the
decision is taken over the profile, the authority context and the lineage.

In PIC-X \texttt{v0.2.2}, the fixed-value part of this policy can be written as
an execution-contract entry such as \(\mathrm{org}=\mathrm{acme}\). The
relational atoms \(\mathrm{SameOrg}\) and \(\mathrm{SoD}_{\mathrm{prin}}\),
which relate predecessor and successor, need binary atoms, which that version
does not implement.

\begin{thebibliography}{99}

\bibitem{gallo2026poc}
N. Gallo.
\newblock Proof-of-Continuity: A Temporal Model for Authority Propagation in
Distributed Systems and AI Agents.
\newblock arXiv:2607.08906, 2026. \url{https://arxiv.org/abs/2607.08906}.

\bibitem{picprotocol}
PIC Protocol.
\newblock \url{https://pic-protocol.org/}.
\newblock Accessed 24 September 2026.

\bibitem{picx}
PIC Protocol project.
\newblock PIC-X.
\newblock Source code, version \texttt{v0.2.2} (commit \texttt{33b9990}), built
on the \texttt{pic-protocol} crate, version \texttt{0.2.3}.
\newblock \url{https://github.com/pic-protocol/pic-x}.
\newblock Accessed 24 September 2026.

\bibitem{sdjwt}
D. Fett, K. Yasuda, and B. Campbell.
\newblock Selective Disclosure for JSON Web Tokens.
\newblock RFC 9901, IETF, November 2025.

\bibitem{hardy1988confused}
N. Hardy.
\newblock The Confused Deputy: Or Why Capabilities Might Have Been Invented.
\newblock \emph{ACM SIGOPS Operating Systems Review}, 22(4):36--38, 1988.

\bibitem{dennis1966programming}
J. B. Dennis and E. C. Van Horn.
\newblock Programming Semantics for Multiprogrammed Computations.
\newblock \emph{Communications of the ACM}, 9(3):143--155, 1966.

\bibitem{miller2003capability}
M. S. Miller, K.-P. Yee, and J. S. Shapiro.
\newblock Capability Myths Demolished.
\newblock Technical Report SRL2003-02, Systems Research Laboratory,
Johns Hopkins University, 2003.

\bibitem{abadi2003history}
M. Abadi and C. Fournet.
\newblock Access Control Based on Execution History.
\newblock In \emph{Proceedings of the Network and Distributed System Security
Symposium (NDSS)}, 2003.

\bibitem{birgisson2014macaroons}
A. Birgisson, J. G. Politz, \'U. Erlingsson, A. Taly, M. Vrable, and M. Lentczner.
\newblock Macaroons: Cookies with Contextual Caveats for Decentralized
Authorization in the Cloud.
\newblock In \emph{Proceedings of the Network and Distributed System Security
Symposium (NDSS)}, 2014.

\bibitem{biscuitspec}
Eclipse Biscuit Project.
\newblock Biscuit Authorization Token Specification.
\newblock \url{https://github.com/eclipse-biscuit/biscuit/blob/main/SPECIFICATIONS.md}.

\bibitem{ucanspec}
UCAN Working Group.
\newblock User Controlled Authorization Network (UCAN) Specification.
\newblock \url{https://github.com/ucan-wg/spec}.

\bibitem{asor2026delegation}
R. Asor.
\newblock Verifiable Attenuated Delegation for AI Agent Chains.
\newblock Internet-Draft draft-asor-wimse-agent-delegation-chain-01, IETF,
September 2026. Work in progress.

\bibitem{atakora2026sadp}
H. Atakora.
\newblock Capability-Bound Delegation Chains with Scoped Context Disclosure for
AI Agents.
\newblock Internet-Draft draft-atakora-wimse-sadp-delegation-01, IETF,
September 2026. Work in progress.

\bibitem{hu2014abac}
V. C. Hu, D. Ferraiolo, R. Kuhn, A. Schnitzer, K. Sandlin, R. Miller, and
K. Scarfone.
\newblock Guide to Attribute Based Access Control (ABAC) Definition and
Considerations.
\newblock NIST Special Publication 800-162, 2014.

\bibitem{ferraiolo2001rbac}
D. F. Ferraiolo, R. Sandhu, S. Gavrila, D. R. Kuhn, and R. Chandramouli.
\newblock Proposed NIST Standard for Role-Based Access Control.
\newblock \emph{ACM Transactions on Information and System Security},
4(3):224--274, 2001.

\bibitem{simon1997sod}
R. T. Simon and M. E. Zurko.
\newblock Separation of Duty in Role-Based Environments.
\newblock In \emph{Proceedings of the 10th IEEE Computer Security Foundations
Workshop (CSFW)}, 1997.

\bibitem{appel1999pca}
A. W. Appel and E. W. Felten.
\newblock Proof-Carrying Authentication.
\newblock In \emph{Proceedings of the 6th ACM Conference on Computer and
Communications Security (CCS)}, 1999.

\bibitem{bauer2005distributed}
L. Bauer, S. Garriss, and M. K. Reiter.
\newblock Distributed Proving in Access-Control Systems.
\newblock In \emph{Proceedings of the IEEE Symposium on Security and Privacy},
2005.

\bibitem{myers1997decentralized}
A. C. Myers and B. Liskov.
\newblock A Decentralized Model for Information Flow Control.
\newblock In \emph{Proceedings of the ACM Symposium on Operating Systems
Principles (SOSP)}, 1997.

\end{thebibliography}

\end{document}
